% ARGUS morphing wing: restructured all-geometry report (v2).
% Sections are fragments in sections/, figures in fig/.
% Compile from this directory: latexmk -pdf argus_rans_validation_report.tex
% Preamble is main.tex's, reused unchanged so the two documents compare like for like.
\documentclass[10pt,a4paper]{article}

\usepackage[a4paper,margin=20mm,top=18mm,bottom=20mm]{geometry}
\usepackage{booktabs}
\usepackage{graphicx}
\usepackage{siunitx}
\usepackage{amsmath}
\usepackage{caption}
\usepackage{subcaption}
\usepackage{longtable}
\usepackage[hidelinks]{hyperref}
\usepackage{microtype}
\usepackage{enumitem}
\usepackage{titlesec}
\usepackage{tikz}
\usetikzlibrary{positioning,arrows.meta}
\titleformat{\section}{\normalfont\large\bfseries}{\thesection}{0.6em}{}
\titleformat{\subsection}{\normalfont\normalsize\bfseries}{\thesubsection}{0.6em}{}
\titleformat{\subsubsection}{\normalfont\normalsize\itshape\bfseries}{\thesubsubsection}{0.6em}{}
\titlespacing*{\subsubsection}{0pt}{1.1ex plus 0.3ex minus 0.2ex}{0.5ex plus 0.2ex}
\titlespacing*{\section}{0pt}{1.7ex plus 0.4ex minus 0.2ex}{0.9ex plus 0.2ex}
\titlespacing*{\subsection}{0pt}{1.3ex plus 0.3ex minus 0.2ex}{0.6ex plus 0.2ex}
\titlespacing*{\paragraph}{0pt}{0.8ex plus 0.2ex minus 0.1ex}{1em}
\setlist[enumerate]{itemsep=1pt,parsep=0pt,topsep=3pt,partopsep=0pt}
\setlist[itemize]{itemsep=1pt,parsep=0pt,topsep=3pt,partopsep=0pt}

\graphicspath{{fig/}}

% Float packing. These change layout only: no content is added or removed by them.
\renewcommand{\topfraction}{0.92}
\renewcommand{\bottomfraction}{0.65}
\renewcommand{\textfraction}{0.06}
\renewcommand{\floatpagefraction}{0.80}
\setcounter{topnumber}{3}
\setcounter{bottomnumber}{2}
\setcounter{totalnumber}{5}
\setlength{\textfloatsep}{7pt plus 2pt minus 2pt}
\setlength{\intextsep}{7pt plus 2pt minus 2pt}
\setlength{\floatsep}{7pt plus 2pt minus 2pt}
\setlength{\abovecaptionskip}{4pt}
\setlength{\belowcaptionskip}{2pt}

\hypersetup{
  pdftitle={ARGUS Morphing Wing: High-Fidelity RANS Validation of the VLM Design Study},
  pdfauthor={Tyler Buchanan},
  pdfsubject={Viscous RANS validation of morphing-wing candidates against VSPAERO/DSO}
}

\captionsetup{font=small,labelfont=bf}
\renewcommand{\arraystretch}{1.0}
\setcounter{tocdepth}{1}
% group-digits=integer: without it siunitx groups the DECIMAL digits too, so
% Sref 1.24092 printed as "1.240,92", which reads as a comma decimal. 25 lines of
% the rendered PDF carried that, including two on page 1. Raised by a peer audit.
\sisetup{group-separator={,},group-minimum-digits=4,group-digits=integer}

\title{ARGUS Morphing Wing: High-Fidelity RANS Validation\\of the VLM Design Study}
\author{Tyler Buchanan}
\date{9 October 2026}

% Generated table rows (sections/gen/, scripts/make_cruise_tables.py). \input is not expandable
% inside an alignment and breaks the booktabs rule that follows it; the TeX primitive is.
\newcommand*{\inputrows}[1]{\csname @@input\endcsname #1 }

\begin{document}

\maketitle

\begin{abstract}
\small\noindent
Six wing surfaces were meshed and solved with steady RANS. One is the rigid baseline. Five are
morphing candidates a VSPAERO vortex-lattice study selected on induced drag alone. Every geometry
is trimmed to its own condition's target lift. All coefficients are on the DSO basis,
$S_\mathrm{ref} = \num{1.24092}$~m$^2$ and $C_\mathrm{ref} = \num{0.39396}$~m. Ninety cases
converged and all eighteen trims are solved. At
Condition~CR, $M = 0.10$ and $C_L = 0.428277635108$, the candidates span $-0.868$ to $+0.232$
counts against a baseline of $188.031$ counts, four of the five below it. At early cruise,
$M = 0.78$ and $C_L = 0.529297087450$, W saves $0.769$ counts against a baseline of $177.839$, F
and H cost $1.796$ and $2.244$, and C and M cost $12.871$ and $14.554$. Late cruise, $M = 0.78$
and $C_L = 0.544114800210$, gives the same order, W at $-0.781$ to M at $+13.728$ against
$187.959$, and W is the cheapest geometry at all three operating points. The cost follows the
aft shock: C and M carry a strong shock at $\eta = 0.80$ and separate behind it, F and H a
moderate one and stay attached, W none. The far-field induced-drag metric the study now ranks on reproduces
the measured Condition~CR ordering to nine pairs of ten; the near-field metric the June selection
was made on reproduces four of ten, and neither resolves the size of the effect. Every candidate
raises the hinge moment about the morph line, by $0.4$ to $6.3$ per cent, and the root-bending
constraint that selected C is not satisfied once computed viscously: C moves from $+6.796$ per
cent in the vortex-lattice study, inside the $6.8$ per cent screen, to $+7.170$ per cent in RANS,
outside it. The hinge moments tabulated for the hinged and twist concepts are common-basis
comparisons rather than the load their own mechanisms carry, so the two morphing concepts cannot
yet be ranked on actuator load.
\end{abstract}

\tableofcontents

\section*{Nomenclature}
\addcontentsline{toc}{section}{Nomenclature}
\label{sec:nomenclature}

Reference quantities are dual and are never mixed. The DSO basis is primary: all three fixed-lift
targets are defined on it and every RANS coefficient here is on it, measured from each case's own
\texttt{forceCoeffs} header rather than declared. The TP-1580 tunnel basis is reported alongside.
The two areas differ by $13.3572/12.0 = \num{1.113100}$, i.e.\ $11.31$~per~cent, against a
morphing effect of order $0.4$~per~cent, so a missed conversion is about 28 times the signal. The
exact reciprocal $\num{0.898392}$ carries a TP-1580 coefficient onto the DSO basis, and these two
are the pair used throughout; $\num{1.113132}$ is the same ratio formed from the rounded metric
areas and is not the reciprocal of $\num{0.898392}$.

{\footnotesize
\begin{longtable}{@{}p{0.15\textwidth}p{0.81\textwidth}@{}}
\toprule
Symbol & Definition, value where fixed, and frame \\
\midrule
\endfirsthead
\toprule
Symbol & Definition, value where fixed, and frame \\
\midrule
\endhead
\bottomrule
\endlastfoot

$S_\mathrm{ref}$, $C_\mathrm{ref}$, $c_\mathrm{ref}$ & Full-wing reference area and chord: DSO
  $\num{1.24092}$~m$^2$ ($13.3572$~ft$^2$) with $C_\mathrm{ref}=\num{0.39396}$~m; TP-1580
  $\num{1.1148}$~m$^2$ ($12.0$~ft$^2$) with $c_\mathrm{ref}=\num{0.3404}$~m, the chord used for
  Reynolds matching to the tunnel \\
$A_\mathrm{ref}$, $l_\mathrm{ref}$ & Half-wing reference area carried by every converged case,
  $\num{0.620462}$~m$^2$ ($2A_\mathrm{ref}$ reproduces the DSO $S_\mathrm{ref}$ to
  $3.22\times10^{-6}$ relative), and the mean aerodynamic chord $\num{0.393957}$~m \\
AR, $b_\mathrm{ref}$ & Aspect ratio $\num{10.7807}$ on the DSO area and $\num{12.0000}$ on
  TP-1580's; reference span $\num{3.6576}$~m ($12.0$~ft), common to both \\
\addlinespace
$C_L$, $C_D$ & Lift and total drag coefficient, each a 200-sample gate window mean, never a last
  sample, on the DSO half-model basis, patch \texttt{wing} \\
$C_M$ & Pitching-moment coefficient about $(1.34,0,0)$~m, nose-up positive; written by
  \texttt{forceCoeffs} and reported per case, with the trim set carrying forces only
  (Table~\ref{tab:qv:qoi}) \\
$L/D$, ct & $C_L/C_D$, both window means, the reference area cancelling; drag count,
  $1\times10^{-4}$ in a force coefficient. A $C_L$ residual in counts is written ``counts of
  $C_L$'' \\
$\Delta C_D$ & Candidate-minus-baseline drag increment at fixed $C_L$, within one condition only \\
$C_{D,p}$, $C_{D,v}$ & Pressure and viscous drag split; closes at Condition~CR and, once $q$
  carries $\rho$, at cruise \\
$C_{Di}$, $C_{Diw}$ & Induced drag, near field (surface integration) and far field (wake plane),
  both VSPAERO columns. The DSO reranked onto $C_{Diw}$ on 19~August 2026. The RANS Trefftz-plane
  induced drag, and the span efficiency it implies, are in Section~\ref{sec:cr:spaneff} \\
$C_{Dtot}$, $C_{Do}$ & VSPAERO total $C_{Do}+C_{Di}$, held apart from any RANS total; and its
  flat-plate skin-friction correlation, whose $Re=1\times10^{7}$ is a VSPAERO placeholder
  entering only that correlation, the VLM itself being inviscid \\
$e$, $E$, $E_w$ & Span efficiency in $C_{Di}=C_L^2/(\pi\,\mathrm{AR}\,e)$; $E$ near field and
  $E_w$ far field as printed by VSPAERO \\
\addlinespace
$c_p$, $c_p^{*}$ & Pressure coefficient $(p-p_\infty)/q$, formed in the solver so the
  normalisation travels with the file; and the sonic value $-0.4931$, derived at $M=0.780302$
  with $\gamma=1.4$. These carry no reference area \\
$c_f$, $c_{f,s}$ & Skin friction $|\tau_w|/q$, and its signed freestream-aligned form
  $m(\tau_w\!\cdot\!\hat{d})/q$ with $m=-1$ calibrated per case, where $\tau_s<0$ marks reversed
  flow \\
$c_l$, $c_l c/C_\mathrm{ref}$ & Sectional lift coefficient (pressure only) and sectional load,
  the spanwise loading quantity \\
$m'_h$, $c_{mh}$ & Sectional hinge moment per unit span, N\,m/m, and its coefficient
  $m'_h/(q_\infty c(y)^2)$ formed on the \emph{local} section chord rather than a reference
  chord. Taken about the swept $x_h/c = 0.62$ line, projected on that line's \emph{local
  tangent}, over aft surfaces only, positive trailing-edge down \\
$M_x$ & Half-wing root bending moment about the root chord line at $y = 0$, N\,m, taken from
  the solver's own moment output rather than integrated from a surface. Dimensional, so it is
  \emph{not} comparable across dynamic pressures; the comparison against the low-order method is
  made on the percentage increase over each method's own baseline, which cancels $q$
  (Section~\ref{sec:rbm}) \\
$M_h$, $C_{Mh,\mathrm{region}}$ & Integrated half-wing hinge moment over
  $0.60 \leq \eta \leq 1.00$, N\,m, and its coefficient
  $M_h/[q_\infty\!\int\! c(y)^2\,\mathrm{d}y]$ over the same region, which avoids selecting one
  of this wing's several defensible reference chords. $M_h$ is dimensional and therefore
  \emph{not} comparable across the three dynamic pressures; $C_{Mh,\mathrm{region}}$ is \\
$C_f$ & Flat-plate correlation $0.026\,Re^{-1/7}$, used only in the $y^+$ stack design and
  distinct from $c_f$ \\
\addlinespace
$M$, $\alpha$, $\alpha_{0L}$ & Mach number; incidence in degrees read from the solver's own
  drag-direction header, never from the case name; zero-lift angle from the bracket fit \\
$U_\infty$, $\rho_\infty$, $p_\infty$, $T_\infty$, $q$ & Freestream speed, density, static
  pressure, temperature and dynamic pressure; $q$ is kinematic at Condition~CR
  ($\tfrac{1}{2}|U|^2=578$~m$^2$/s$^2$) and absolute at cruise \\
$\gamma$ & Ratio of specific heats: $1.4$ for $c_p^{*}$, and the file's own $\num{1.4010834}$ for
  the late-cruise dynamic-pressure closure \\
$\nu$, $\mu$, $\nu_t$, $\tilde{\nu}$ & Kinematic and dynamic viscosity (the cruise $\mu$ is the
  model-scale Reynolds-matched value, not the physical one); eddy viscosity $\nu_t=k/\omega$;
  Spalart-Allmaras transported variable \\
$k$, $\omega$, $Tu$ & Turbulence kinetic energy, specific dissipation rate, and freestream
  turbulence intensity $\sqrt{2k/3}/U_\infty$ \\
$Re_{C_\mathrm{ref}}$, $Re_{c_\mathrm{ref}}$ & Reynolds number on the DSO and on the TP-1580 chord
  respectively; the subscript names which chord, and both are quoted for every condition \\
$\tau_w$, $\tau_s$, $u_\tau$, $y^+$ & Wall shear stress, its signed streamwise component, friction
  velocity, and $y\,u_\tau/\nu$, quoted throughout on the \emph{cell-centre} convention the solver
  evaluates, on which a stack designed at the full first-cell height reads half its nominal
  target (Section~\ref{sec:setup:mesh}) \\
\addlinespace
$\eta$, $x/c$, $x_h/c$ & Spanwise station on $b_\mathrm{ref}/2$ unless a table names the extractor
  axis, which is $0.19$~per~cent inboard of it; chordwise fraction of local chord; morph start or
  hinge line, $0.62$ continuous and $0.70$ hinged \\
$\xi$, $S(\xi)$, $A/c$, $\Delta z_c$, $\Delta\theta$ & Morph coordinate
  $(x/c-x_h/c)/(1-x_h/c)$; smoothstep $3\xi^2-2\xi^3$; commanded camber amplitude on the local
  chord, positive trailing-edge down; camber displacement $-A(\eta)S(\xi)$ applied equally to both
  surfaces so thickness is preserved; commanded twist in degrees \\
\addlinespace
$W$, $n$, $\tau$ & Convergence-gate window (200 force samples); converged bracket points in a fit;
  Kendall rank correlation (wall shear is $\tau_w$) \\
B, C, F, M, W, H & The six geometries: rigid baseline; continuous trailing-edge camber (C, F, M);
  distributed twist (W); conventional hinged (H) \\
B$^{*}$ & The corrected-interpolation baseline the DSO ranks against, which no RANS case has ever
  meshed \\
Condition~CR & The $M = 0.10$ low-speed state the design study selected its candidates at.
  \textbf{CR does not stand for cruise}: the two cruise states are at $M = 0.78$ and are always
  written ``early cruise'' and ``late cruise''. Unqualified ``cruise'' in this report means those
  two and never Condition~CR \\
\texttt{SW*}, \texttt{CMP*}, \texttt{LC*} & Case prefixes at Condition~CR, early cruise and late
  cruise respectively \\
DSO, VLM, RANS, EET & Design-space optimisation, the VSPAERO vortex-lattice study that selected
  the candidates; vortex-lattice method; Reynolds-averaged Navier-Stokes; Energy Efficient
  Transport, the NASA TP-1580 high-aspect-ratio supercritical-wing model \\
\end{longtable}
}

\section{Objectives and the question}
\label{sec:intro}
\label{sec:objectives}
\textbf{Purpose.} ARGUS asks how the outer wing of a transport aircraft should adapt during
cruise. A low-order VSPAERO study searched that design space and selected the candidate
shapes~\cite{zheng}. The search is inviscid, wing-only, and ranks on induced drag. This work is
the viscous audit of that search.

\textbf{Study hierarchy.} The design study evaluates hundreds of geometries cheaply. This
campaign evaluates six of them with steady RANS at the matched Reynolds number. The design study
proposes candidates. The RANS says what they cost at trimmed lift.

\textbf{The job.} Six wing surfaces, one rigid baseline and five morphing candidates, were meshed and
solved with steady RANS at three operating points, every geometry trimmed to that point's target
lift. The low-order study selected the candidates at the low-speed point, Condition~CR, so that is
where it is tested against the RANS; the two cruise points are the results the low-order study
does not have.
\textbf{Question.} This report answers two. First, does the low-order method order the candidates
the way the RANS orders them? Second, what does the RANS say about the candidates themselves?

\textbf{Reading order.} The report follows the order of the work. Section~\ref{sec:condgeom} states
the operating points and the six geometries. Section~\ref{sec:method} states the OpenFOAM setup,
the automated pipeline, the trim convention and the convergence gate, each once.
Section~\ref{sec:qoi} lists every quantity of interest, its frame, and whether the low-order study
produces it. Section~\ref{sec:results} reports each quantity in turn, RANS beside VLM wherever a
VLM value exists. Section~\ref{sec:lo} answers whether the VLM can drive the design study, and
Section~\ref{sec:concl} states what the whole of it establishes and what it does not.
Appendices~\ref{sec:app:sections} and~\ref{sec:deltamaps} carry the surface pressure and skin
friction, by station and over the whole wing.

\section{Conditions and geometries}
\label{sec:condgeom}
\label{sec:geom}

\textbf{Purpose.} This section fixes the three states every number in this report is read at, and
says what was changed on the wing. Five candidates are compared against one rigid baseline. Each
moves a different part of the outer wing by a different amount. Every geometric number here is
read from the design files and the delivered section schedules, and mesh statistics are quoted
where they bear on the geometry.
\subsection{Operating points}
\label{sec:cond}
\label{sec:cond:points}

\textbf{Purpose.} Three states are fixed here, and every number in this report is read at one of
them. Two are cruise states of the ARGUS aircraft at $M = 0.78$. One is the low-speed point the DSO
optimisation selected its candidates at. Table~\ref{tab:cond:points} is the record of what ran.
\textbf{``Condition~CR'' is the low-speed point, and CR does not stand for cruise.} It is the
$M = 0.10$ state the design study selected its candidates at, and it is named Condition~CR
throughout this report and in the delivered packages. The two states that \emph{are} cruise are
always written ``early cruise'' and ``late cruise'', both at $M = 0.78$. The distinction is
load-bearing and the wrong reading inverts results: where Section~\ref{sec:hinge} reports a
penalty that is smaller at Condition~CR than at early cruise, reading CR as cruise makes the
sentence contradict itself. Wherever this report says ``cruise'' without qualification it means
the two $M = 0.78$ states and never Condition~CR.

\textbf{Method.} Each cruise state is pinned by mass, altitude and Mach number. Standard-atmosphere
density and speed of sound follow from the altitude. The target lift coefficient is then
$C_{L,\mathrm{target}} = W / (q_\infty S_\mathrm{ref})$. Every value in
Table~\ref{tab:cond:points} is read from the recipes' \texttt{CONDITION.json} and from each built
case's own \texttt{forceCoeffs} block.

\textbf{Scale.} The cruise states are defined at aircraft scale, on a chord of $3.196$~m. The
meshed geometry is model scale, on $C_\mathrm{ref} = \num{0.393957}$~m. Viscosity is set so that
Reynolds number on the model chord equals the aircraft-scale value. The table therefore carries a
model-scale $\mu$ beside an aircraft Reynolds number.

\begin{table}[htbp]
  \centering
  \footnotesize
  \caption{The three operating points. Coefficients are on the DSO basis,
  $S_\mathrm{ref} = \num{1.24092}$~m$^2$ and $C_\mathrm{ref} = \num{0.39396}$~m; the RANS cases
  solve a half wing at $A_\mathrm{ref} = \num{0.620462}$~m$^2$. Reynolds number is on the DSO
  chord: at Condition~CR it is $0.79$ million on the TP-1580 chord of $\num{0.3404}$~m, which is
  the figure in wider circulation. Condition~CR is solved incompressible, so its pressure is
  kinematic and its temperature is not defined; the tabulated $q_\infty$ is
  $\tfrac{1}{2}\rho_\infty|U_\infty|^2$ on the $\rho_\infty = 1.225$ its \texttt{forceCoeffs}
  block carries. Our cruise $q_\infty$ runs $0.303$~per~cent above Zheng's
  Table~18~\cite{zheng} at both states, common mode within our runs.}
  \label{tab:cond:points}
  \begin{tabular}{lrrr}
    \toprule
                                        & Condition CR      & Early cruise            & Late cruise \\
    \midrule
    Case prefix                         & \texttt{SW*}      & \texttt{CMP*}           & \texttt{LC*} \\
    Flight state                        & DSO-matched       & \num{74500}~kg, 10~km   & \num{56000}~kg, 12~km \\
    Mach                                & $0.10$            & $0.78$                  & $0.78$ \\
    Target $C_L$                        & $0.428277635108$  & $0.529297087450$        & $0.544114800210$ \\
    $|U_\infty|$ (m/s)                  & $34.0$            & $233.934028$            & $230.501789$ \\
    $\rho_\infty$ (kg/m$^3$)            & $1.225$           & $0.41271$               & $0.31083$ \\
    $T_\infty$ (K)                      & n/a               & $223.15$                & $216.65$ \\
    $p_\infty$ (Pa)                     & n/a (gauge)       & $\num{26495.90}$        & $\num{19373.96}$ \\
    $q_\infty$ (Pa)                     & $708.05$          & $\num{11292.80}$        & $\num{8257.37}$ \\
    Viscosity                           & $\nu\,1.46\times10^{-5}$ & $\mu\,1.7961143\times10^{-6}$ & $\mu\,1.7523444\times10^{-6}$ \\
    $Re_{C_\mathrm{ref}}$               & $0.92\times10^{6}$ & $2.1176454\times10^{7}$ & $1.6107442\times10^{7}$ \\
    Solver                              & incompressible    & compressible            & compressible \\
    Turbulence model                    & kOmegaSST         & Spalart--Allmaras       & Spalart--Allmaras \\
    Wall treatment                      & resolved          & modelled                & modelled \\
    \bottomrule
  \end{tabular}
\end{table}
\textbf{Why the two cruise states differ.} The aircraft burns fuel and steps from 10 to 12~km, so
mass falls by a quarter and density falls with the altitude. Dynamic pressure falls faster than
weight, so the required lift coefficient rises by $2.80$~per~cent while Reynolds number drops by
$24$~per~cent.

\textbf{The two states close against each other.} At fixed reference area the ratio of targets
equals the mass ratio over the dynamic-pressure ratio. The targets give $1.0279951$. Mass and
dynamic pressure give $1.0279974$, agreeing to $2.3\times10^{-6}$ relative. The two points are one
flight state read at two masses.

\subsection{The wing every candidate starts from}
\label{sec:geom:baseline}

\textbf{Baseline.} All six surfaces are the same NASA EET AR-12 cruise wing. Only the outer wing
differs between them. Fuselage, tail and high-lift system are not modelled.
Table~\ref{tab:geom:planform} gives the planform an experimentalist would need.

\textbf{Section.} The aerofoil is a supercritical section of $11.97$~per~cent thickness. Its
trailing edge is blunt at $0.63$~per~cent of local chord. That edge is $1.03$~mm at the tip and
sets the surface cell size everywhere.

\begin{table}[htbp]
  \centering
  \footnotesize
  \caption{Baseline planform, read from the geometry sidecar of the surface the cases mesh,
  \texttt{geometry/derived/wing3d/baseline\_oml\_placed.json}. Measured values are taken from the
  lofted surface; targets are the registered references. The half-wing planform area integrates
  to \num{0.620463}~m$^2$ against the $A_\mathrm{ref} = \num{0.620462}$~m$^2$ the cases carry.
  The meshing surface adds a $20$~mm root extension through the symmetry plane as scaffolding,
  which is cut away and is not part of the wing.}
  \label{tab:geom:planform}
  \begin{tabular}{lrrl}
    \toprule
    Quantity & Measured & Target & Note \\
    \midrule
    Root chord (m)                  & $0.639125$  & $0.640081$ & \\
    Tip chord (m)                   & $0.152172$  & $0.152400$ & taper ratio $0.238$ \\
    Semispan (m)                    & \multicolumn{2}{r}{$b_\mathrm{ref}/2 = 1.8288$} & half model on a symmetry plane \\
    Leading-edge sweep (deg)        & $28.8447$   & $28.8620$  & quarter chord $27$ \\
    Tip minus root twist (deg)      & $-4.129406$ & $-4.11495$ & washout \\
    Root incidence (deg)            & $1.852649$  & $1.873$    & placement transform \\
    Dihedral rise over semispan (m) & $0.130498$  & $0.139$    & placement transform \\
    Half planform area (m$^2$)      & \multicolumn{2}{r}{$0.620463$} & integrated from the section table \\
    Thickness, trailing edge        & \multicolumn{2}{r}{$11.97\,\%$ chord; TE $0.63\,\%$ local chord} & TE $1.03$~mm minimum \\
    \bottomrule
  \end{tabular}
\end{table}

\textbf{Span coordinate.} Stations below are quoted in $\eta = y/(b/2)$. The wing is lofted
through 19 defining stations. Those stations carry the morphing, and every realised value in this
section is read at them.

\subsection{Two lineages of the same wing}
\label{sec:geom:lineage}

\textbf{Why there are two.} The delivered design file lofts the wing through 19 stations but
carries only 9 distinct aerofoils: ten stations inherit a neighbour's section instead of a
blended one, so thickness is held constant across runs of stations and then steps. A corrected
file was issued afterwards in which all 19 stations carry their own aerofoil. The meshed
baseline is the uncorrected file. All five candidates are built on the corrected one.

\textbf{What the difference is.} Figure~\ref{fig:geom:lineage} shows it. Both wings are the
same planform at the same twist, so the difference is thickness and nothing else. They agree
exactly at every station both files define and differ only between those stations, which is
why the uncorrected thickness distribution is a staircase and the corrected one a smooth
taper. The difference is zero at and inboard of $\eta = 0.450$. The first station that differs,
$\eta = 0.515$, is also the largest, at $0.407$ percentage points of $t/c$ and $0.209$~per~cent
of chord in surface ordinate. Camber differs by at most $0.0034$~per~cent of chord, two orders
of magnitude below the thickness difference.

\textbf{What it costs the comparison.} Nothing for the candidate ranking. All five candidates
share the corrected lineage, so the difference is common to both sides of every
candidate-to-candidate number in this report and cancels from it. It enters only where a
candidate is differenced against the meshed baseline, and there it is bounded by the figures
above.

\begin{figure}[htbp]
  \centering
  \includegraphics[width=0.94\linewidth]{fig/fig_lineage_staircase.pdf}
  \caption{The two lineages of the baseline wing. Left: maximum $t/c$ against $\eta$, the
  uncorrected file (9 distinct aerofoils across 19 stations) against the corrected file (19).
  The uncorrected curve holds thickness constant across runs of stations and steps between
  them; the corrected curve tapers. Centre: the two adjacent uncorrected stations either side
  of the largest step, overlaid, which is what an inherited rather than blended section looks
  like. Right: the surface ordinate difference $|\Delta z|/c$ across that step for each
  lineage; the uncorrected step reaches $0.25$~per~cent of chord where the corrected reaches
  $0.06$. Measured from the two design files by
  \texttt{scripts/measure\_lineage\_pair.py} into \texttt{data/lineage\_measured.json}.}
  \label{fig:geom:lineage}
\end{figure}

\subsection{How the morphing is commanded}
\label{sec:geom:law}

Three mechanisms appear across the five candidates. Each acts on a different part of the section.

\textbf{Continuous trailing-edge camber (C, F, M).} Camber changes behind a fixed start line at
$x_h/c = 0.62$. Let
\begin{equation}
  \xi = \frac{x/c - x_h/c}{1 - x_h/c}, \qquad 0 \le \xi \le 1 .
\end{equation}
The camber displacement follows a smoothstep in $\xi$,
\begin{equation}
  S(\xi) = 3\xi^{2} - 2\xi^{3}, \qquad
  \Delta z_c(x,\eta) = -A(\eta)\,S(\xi), \qquad x/c \ge x_h/c .
\end{equation}
Both $S$ and its slope vanish at the start line. The moving surface therefore leaves the fixed
forward aerofoil without a step or a slope break. Upper and lower surfaces are displaced by the
same increment, so thickness is preserved by construction. The delivered sections carry zero
change in $t/c$ and a trailing-edge displacement equal to $-A/c$ to $10^{-9}$. Positive $A$ is
trailing edge down. The amplitude $A/c$ is normalised on the \emph{local} chord.

\textbf{Distributed twist (W).} The whole local section rotates about $x/c = 0.25$. The aerofoil
profile is unchanged and only incidence moves. The command is an incremental twist in degrees,
added to the baseline washout. This is the one candidate whose leading edge moves.

\textbf{Conventional hinged (H).} Two aileron-like segments rotate rigidly about a hinge at
$x_h/c = 0.70$. Deflection is constant within a segment and steps between them. The slope breaks
at the hinge, which is what separates this concept from camber morphing. It is modelled as a
zero-gap outer mould line, with no gap, seal or bracket represented.

\textbf{Command layout.} Each continuous concept carries five spanwise commands. These are design
variables and not geometric sections. Values between them are interpolated onto the 19 defining
stations before the surface is lofted. Table~\ref{tab:geom:command} gives the commands as the
design files hold them.

\begin{table}[htbp]
  \centering
  \scriptsize
  \setlength{\tabcolsep}{4pt}
  \caption{The six geometries and the command each flies, read from the DSO design files
  (\texttt{<case>\_design.json}) and \texttt{registry/candidates.yaml}. Camber commands are $A/c$
  on the local chord at $\eta = 0.60$, $0.70$, $0.80$, $0.90$, $1.00$. Twist commands are
  $\Delta\theta$ in degrees at $\eta = 0.68$, $0.76$, $0.84$, $0.92$, $1.00$, positive leading
  edge up. The hinged command is the inboard and outboard segment deflection in degrees, positive
  trailing edge down. The STL column is the SHA-256 prefix of the placed surface the cases mesh.
  All five candidates were selected at Condition~CR, $M = 0.10$, $C_L = 0.428277635108$
  (Table~\ref{tab:cond:points}).}
  \label{tab:geom:command}
  \begin{tabular}{lllll}
    \toprule
    ID & DSO case & Concept & Command & STL \\
    \midrule
    B & \texttt{baseline\_wing\_only\_refined} & rigid     & none & \texttt{20bd6acd} \\
    C & \texttt{cte\_i002\_c04}  & TE camber & $0.016427$, $0.021826$, $0.033726$, $0.015955$, $0.000015$ & \texttt{790c3ded} \\
    F & \texttt{cfft\_b02\_c01}  & TE camber & $0.012607$, $0.013869$, $0.012069$, $0.012284$, $0.000461$ & \texttt{82bf2e1e} \\
    M & \texttt{mcv2\_i002\_c01} & TE camber & $0.008027$, $0.025978$, $0.034968$, $0.018199$, $0.000218$ & \texttt{dae22d0c} \\
    W & \texttt{cffw\_b01\_c01}  & twist     & $+1.7997$, $+1.1260$, $+1.2782$, $+0.2573$, $-0.9994$ & \texttt{200fa64c} \\
    H & \texttt{chc\_g02\_c06}   & hinged    & $2.4954$ inboard, $1.6636$ outboard & \texttt{18329aa3} \\
    \bottomrule
  \end{tabular}
\end{table}

\textbf{Two generations, not five designs.} C and M come from the July selection. F, W and H come
from the August reselection, after the design study moved from near-field to far-field induced
drag. C and F are the same concept ranked under two metrics. F's peak amplitude is $41$~per~cent
of C's.

\subsection{What each candidate actually moves}
\label{sec:geom:delta}

\textbf{Method.} Commands are interpolated onto the defining stations, so the realised peak sits
at a station rather than at a control point. The chord tapers by a factor of $4.2$ from root to
tip. A constant $A/c$ therefore produces a falling physical displacement outboard.
Table~\ref{tab:geom:delta} gives commanded amplitude and realised displacement together, and
Figure~\ref{fig:geom:displacement} shows the same thing on the wing.

\begin{figure}[htbp]
  \centering
  \includegraphics[width=\linewidth]{fig/concept_displacement_3d.png}
  \caption{Where each concept moves the surface, seen from above and below. Panel letters are the
  geometry IDs of Table~\ref{tab:geom:command}. Span runs left to right with the root at left and
  chord up the page. Colour is displacement from the baseline surface measured along the
  baseline's own outward normal, so the same colour means the same physical direction on both
  surfaces: the four camber concepts move the lower surface down and the upper surface up over
  the aft wing, and W, a rigid rotation about $x/c = 0.25$, moves the whole section including the
  leading edge, which is the signature that distinguishes a twist from a camber morph. One
  colour scale spans every panel, $\pm 6.42$~mm at the $99.5$th percentile, because per-panel
  autoscaling would make five different magnitudes look alike. Surfaces are matched by nearest
  point rather than by vertex index, since the three later concepts carry \num{195266} triangles
  against the baseline's \num{192588}. Generated by \texttt{scripts/render\_concept\_3d.py}.}
  \label{fig:geom:displacement}
\end{figure}

\begin{table}[htbp]
  \centering
  \footnotesize
  \setlength{\tabcolsep}{4pt}
  \caption{Realised geometry change at the 19 defining stations. Computed from the delivered
  schedules: \texttt{<case>\_section\_schedule.csv} for C, F, M and H, and
  \texttt{cffw\_b01\_c01\_twist\_schedule.csv} for W. ``Peak command'' is the largest amplitude
  reached at a defining station. It is $A/c$ for the camber candidates and degrees for W and H.
  ``Peak TE'' is the largest trailing-edge displacement in millimetres, with the station it
  occurs at. For W the trailing edge travels on an arc, and the entry is
  $0.75\,c\sin\Delta\theta$ derived from the tabulated twist and chord. Displacement is trailing
  edge down in every case except W outboard of $\eta = 0.95$.}
  \label{tab:geom:delta}
  \begin{tabular}{llrrrrr}
    \toprule
    ID & Acts over & Stations & Peak command & at $\eta$ & Peak TE (mm) & at $\eta$ \\
    \midrule
    C & $\eta$ $0.6125$ to tip & 12 & $0.031465$   & $0.781$ & $6.894$ & $0.781$ \\
    F & $\eta$ $0.6125$ to tip & 12 & $0.013689$   & $0.710$ & $3.452$ & $0.6125$ \\
    M & $\eta$ $0.6125$ to tip & 12 & $0.033260$   & $0.781$ & $7.288$ & $0.781$ \\
    W & $\eta$ $0.6125$ to tip & 12 & $+1.7435$~deg & $0.6775$ & $5.719$ & $0.6775$ \\
    H & $\eta$ $0.710$ to $0.970$ & 8 & $2.4954$~deg & $0.710$ & $3.145$ & $0.710$ \\
    \bottomrule
  \end{tabular}
\end{table}

\textbf{Result.} Every candidate moves the wing by millimetres. The largest movement anywhere in
the set is $7.29$~mm, on M. That is $1.1$~per~cent of the root chord. C follows at $6.89$~mm and
W at $5.72$~mm. F and H are the smallest at $3.45$ and $3.15$~mm.

\textbf{Taper moves the peak.} F shows the effect plainly. Its command peaks at $\eta = 0.710$,
but its physical displacement peaks at $\eta = 0.6125$, the innermost morphed station. The chord
there is $0.270$~m against $0.241$~m. C and M peak in both measures at the same station, because
their commands rise steeply enough to beat the taper.

\textbf{The tip is where the concepts part.} All three camber candidates close to near zero at the
tip: $0.002$~mm on C, $0.033$~mm on M and $0.070$~mm on F. W does the opposite. It reaches
$+1.74$~deg of washin at $\eta = 0.6775$, crosses zero between $\eta = 0.926$ and $0.970$, and
ends at $-0.999$~deg of washout at the tip. W is the only bidirectional candidate in the set, and
the only one that unloads the tip by rotating it.

\textbf{H is discrete where the others are smooth.} Its inboard segment holds $2.4954$~deg from
$\eta = 0.710$ to $0.852$. Its outboard segment holds $1.6636$~deg to $\eta = 0.970$. Deflection
steps by $0.832$~deg at the segment break. The outer $3$~per~cent of span carries no deflection
at all.

\subsection{Where the morphing acts in span}
\label{sec:geom:span}

\textbf{Region.} The continuous concepts act over $\eta = 0.60$ to the tip. That interval is
$25.2$~per~cent of the half-wing planform, taper making it far less than the $40$~per~cent of
span it occupies. Only the chord aft of $x_h/c = 0.62$ moves, so the moving surface is
$9.6$~per~cent of the half-wing planform. H acts over $\eta = 0.710$ to $0.970$, which is
$15.4$~per~cent of planform, and moves $4.6$~per~cent aft of its hinge.

\textbf{Measured effect.} Where that geometry appears in the loading is shown in
Section~\ref{sec:cr:loads}, Figure~\ref{fig:geom:spanwise}: the difference from baseline sits over
the morphed span in every case.

\subsection{Meshed surfaces and their statistics}
\label{sec:geom:mesh}

\textbf{Route.} All six surfaces are lofted by one procedure, step for step, but that procedure
was run twice: B, C and M went through it in July and F, W and H in August, and the two runs
differ only in the tip-cap construction. Each is tessellated at
200 chordwise points per surface and 240 spanwise stations, with the 19 defining stations
retained. The same placement transform and the same blunt-trailing-edge convention follow.
Table~\ref{tab:geom:mesh} gives the surfaces and the meshes built on them.

\begin{table}[htbp]
  \centering
  \footnotesize
  \setlength{\tabcolsep}{5pt}
  \caption{Meshed surfaces and as-built mesh size, on the meshes every result in this report is
  taken from. Triangle counts are from each surface's own sidecar and split by delivery, not by concept.
  Condition~CR counts are the \texttt{checkMesh} figures carried by the run cards; cruise counts
  are from each mesh's own \texttt{polyMesh/owner} header. The surface gate is
  \texttt{surfaceCheck} on the case copy: zero open edges, zero over-connected edges, one zone
  and one normal orientation. Mesh recipe and layer stacks are in Section~\ref{sec:setup:mesh}. Source: rows generated by
  \texttt{scripts/make\_cruise\_tables.py}.}
  \label{tab:geom:mesh}
  \begin{tabular}{lllrrrl}
    \toprule
    ID & Concept & Triangles & Cells, Condition CR & Cells, early cruise & Cells, late cruise & Surface gate \\
    \midrule
\inputrows{sections/gen/tab_geom_mesh_rows.tex}
    \bottomrule
  \end{tabular}
\end{table}

\textbf{What the mesh size tracks.} Nothing but the condition: the six Condition~CR meshes lie
within $0.10$~per~cent of one another and the twelve cruise meshes within $0.15$~per~cent, so
every candidate-minus-baseline difference is taken between meshes of the same size.

\textbf{The integrated surface.} The wing patch carries \num{5276562} to \num{5277100} faces
across the eighteen trim cases, a spread of about one face in $10^{4}$. The surface discretisation is
therefore common to both sides of every candidate-minus-baseline difference. The meshing surface
also carries the $20$~mm root extension, so its STL area of $1.3497$~m$^2$ stands $5.49$~per~cent
above the wing's wetted area of $1.2794$~m$^2$. Wetted area is taken from the solver wing patch.

\subsection{What was changed, in one place}
\label{sec:geom:summary}

\textbf{Decision.} The five candidates are three concepts and two selection generations. They act
on the outer $40$~per~cent of span and move between $3.1$ and $7.3$~mm of trailing edge. C, F and
M differ in amplitude and spanwise placement rather than in mechanism, and all three close to
zero at the tip. W and H are the structurally distinct pair: W rotates whole sections
bidirectionally and is the only candidate that moves a leading edge, and H breaks the surface
slope at a hinge.

\section{OpenFOAM setup and the automated pipeline}
\label{sec:method}

\textbf{Purpose.} Every number in this report depends on a solver recipe, a trim convention and a
convergence gate, and every case reached its number down one automated route. This section fixes
each of them once, so that no result section has to restate them and no two results can rest on
different ones.
\subsection{Solver, boundary conditions and mesh}
\label{sec:setup}
\label{sec:setup:solver}

\textbf{Recipe.} Every value here was read from the recipe folder each case was built from, and
the recipes were confirmed against the built cases \texttt{SWB\_a1p60} and \texttt{CMPB\_a0p87}.
Per-geometry surface and cell counts are in Table~\ref{tab:geom:mesh}.
\textbf{Method.} Both conditions run \texttt{foamRun} under OpenFOAM-org~12~\cite{openfoam},
steady state, SIMPLEC. Condition~CR uses the \texttt{incompressibleFluid} module at
$\nu = \num{1.46e-5}$~m$^2$/s. Cruise uses the \texttt{fluid} module in transonic mode with a
perfect gas. Table~\ref{tab:setup:numerics} gives both setups.

\textbf{Two models, by design.} Condition~CR runs kOmegaSST wall resolved and cruise runs
Spalart-Allmaras wall modelled. The criterion is the drag split: skin friction cancels between
the members of a morph pair, and Spalart-Allmaras is the stronger model on the drag-rise
derivative that does not. \texttt{nutLowReWallFunction} sets $\nu_t = 0$ at the wall and is the
resolved-sublayer condition. Within a condition the model and the wall treatment are uniform.

\textbf{Numerics.} The cruise schemes and solution control are byte-copied from the project's
Onera~M6 case. That case reproduces AGARD AR-138 run 308~\cite{agard138}. Condition~CR runs
limited gradients, a limited surface-normal gradient and one non-orthogonal corrector.

\textbf{Residuals.} Residual targets sit below anything these solves reach. Each run therefore
ends at \texttt{endTime}, and the force gate of Section~\ref{sec:setup:gate} decides. Achieved
residuals are flat rather than falling. Condition~CR reaches $U_x \sim 10^{-9}$ and cruise
$U_x \sim \num{1.8e-7}$.

\begin{table}[htbp]
  \centering
  \scriptsize
  \caption{Discretisation and solution control, read from the two recipe folders and confirmed
  against the built cases. Residual targets are set below anything the solve reaches, so the run
  ends at \texttt{endTime} and the gate of Table~\ref{tab:setup:gate} decides; the achieved
  plateau levels are given alongside. The cruise column serves both $M = 0.78$ states, which
  differ only in freestream state and layer stack.}
  \label{tab:setup:numerics}
  \begin{tabular}{lll}
    \toprule
     & Condition CR & Cruise (early and late) \\
    \midrule
    Module, fluid        & \texttt{incompressibleFluid}, $\nu$ $\num{1.46e-5}$
                         & \texttt{fluid} transonic, $Cp$ $1005$, $Pr$ $0.71$ \\
    Turbulence model     & kOmegaSST & Spalart-Allmaras \\
    $\nu_t$ wall condition & \texttt{nutLowReWallFunction} & \texttt{nutUSpaldingWallFunction} \\
    Gradient             & \multicolumn{2}{l}{Gauss linear, cell limited on $U$ and turbulence} \\
    Convection, $U$      & bounded linearUpwind & linearUpwind \\
    Convection, turbulence & bounded limitedLinear 1 & linearUpwind \\
    Convection, energy   & n/a & \texttt{filteredLinear2 0.2 0} \\
    Laplacian, snGrad    & limited corrected $0.33$ & corrected \\
    $p$ solver           & GAMG $10^{-8}$, relTol $0.01$ & GAMG $10^{-8}$, relTol $0.1$ \\
    $U$, turbulence solver & smoothSolver $10^{-9}$ & GAMG $10^{-8}$, relTol $0.1$ \\
    Non-orth. correctors & 1 & 0 \\
    Relaxation           & $U$ $0.5$, all $0.5$ & $p$ $1$, $U$ $0.9$, $e$ $0.8$ \\
    Residual targets     & $p$ $10^{-9}$, $U$ $10^{-10}$
                         & $p$ $10^{-9}$, $U$ $10^{-10}$, $e$ $10^{-9}$ \\
    Residuals achieved   & $U_x$ $\sim10^{-9}$, $p$ $\sim10^{-8}$
                         & $U_x$ $\sim\num{1.8e-7}$, $p$ $\sim\num{2e-8}$ \\
    \bottomrule
  \end{tabular}
\end{table}

\textbf{Domain.} The domain is free air. One background block spans $60\times30\times60$~m,
which puts the far field $16.4$ semispans out. The root plane is a symmetry plane and the wing
is \texttt{noSlip}. Table~\ref{tab:setup:bc} gives the patch conditions per field.
\textbf{Incidence.} Incidence enters through the freestream vector and the force axes.
Condition~CR runs $U_\infty = (33.991249,\,0,\,0.771369)$~m/s. Early cruise runs
$(233.907340,\,0,\,3.533537)$~m/s. Each case's $\alpha$ is read back from the \texttt{dragDir}
header the solver wrote, never from the case name.

\textbf{Turbulence inflow.} Every case is fully turbulent from the leading edge. Condition~CR
carries $k = \num{1.734e-3}$~m$^2$/s$^2$ and $\omega = 14.25$~s$^{-1}$. That is a turbulence
intensity of $0.10$~per~cent and $\nu_t/\nu = 8.3$. Cruise carries
$\tilde{\nu} = \num{1.3056e-5}$~m$^2$/s, and zero at the wall.

\begin{table}[htbp]
  \centering
  \scriptsize
  \caption{Boundary conditions as built, per patch and field. Condition~CR fixes pressure at the
  outlet and leaves the other far-field patches \texttt{zeroGradient}; the compressible cases use
  the \texttt{freestream} pair on all four far-field patches. Forces are integrated over the
  \texttt{wing} patch alone.}
  \label{tab:setup:bc}
  \begin{tabular}{lll}
    \toprule
    Patch & Condition CR & Cruise (early and late) \\
    \midrule
    \texttt{wing}   & $U$ \texttt{noSlip}, $p$ \texttt{zeroGradient}
                    & $U$ \texttt{noSlip}, $p$ \texttt{zeroGradient} \\
                    & $k$ \texttt{kLowReWallFunction}
                    & $T$ \texttt{zeroGradient}, $\tilde{\nu}$ fixed $0$ \\
                    & $\omega$ \texttt{omegaWallFunction}
                    & $\nu_t$ \texttt{nutUSpaldingWallFunction} \\
                    & $\nu_t$ \texttt{nutLowReWallFunction}
                    & $\alpha_t$ \texttt{alphatWallFunction} \\
    \midrule
    \texttt{symmetry} & \texttt{symmetry} & \texttt{symmetryPlane} \\
    \midrule
    \texttt{inlet}, \texttt{topBottom}, & $U$, $k$, $\omega$ \texttt{inletOutlet}
                    & $U$ \texttt{freestreamVelocity}, $p$ \texttt{freestreamPressure} \\
    \texttt{outboard} & $p$ \texttt{zeroGradient}, $\nu_t$ \texttt{calculated}
                    & $T$, $\tilde{\nu}$ \texttt{inletOutlet}, $\nu_t$ \texttt{calculated} \\
    \midrule
    \texttt{outlet} & $p$ \texttt{fixedValue} $0$, $U$ \texttt{inletOutlet}
                    & as the other far-field patches \\
    \bottomrule
  \end{tabular}
\end{table}

\label{sec:setup:mesh}
\textbf{Mesh.} One mesh recipe serves both conditions, and only the prism stack differs. The
surface cell is set by the $1.03$~mm blunt trailing edge, not by wall treatment or Reynolds
number. Refinement level 11 gives a $\num{0.488}$~mm surface cell and $2.11$ cells across the
trailing-edge base. The two recipes agree on wall-face count to three faces in
$5.28\times10^{6}$, which is the measured form of the claim that one recipe serves both
conditions. Table~\ref{tab:setup:mesh} gives the recipe and what was built from it.
\textbf{Layers.} Condition~CR asks 11 layers at a $4.5$~$\mu$m first cell and builds $9.79$ of
them. Early cruise asks 5 at $43.7$~$\mu$m and builds $4.96$ on every cruise mesh; late cruise
asks 5 at $56.1$~$\mu$m and builds $4.97$. Achieved stack heights are $95.2$, $99.5$ and
$99.6$~per~cent of the request, read from each mesh's own log; an achieved stack can only fall
below the request, and every one does. Extrusion covers $97.1$~per~cent of wall faces at
Condition~CR and $99.91$ to $99.93$~per~cent at cruise.

\textbf{Wake boxes and cell cap.} Both conditions carry two Trefftz wake boxes over the aft
wing and near wake, at levels 7 and 6, under a $130$~million cell cap. Every one of the twelve fails three to five \texttt{checkMesh} checks,
drawn from skewness, face tetrahedra, small determinant and concavity, plus a
small-interpolation-weight flag on C's early-cruise mesh alone; maximum non-orthogonality is $72.4^{\circ}$ to
$74.9^{\circ}$.

\textbf{Achieved $y^+$ is the design value.} The stack is sized to put the first-cell top at the
nominal target. OpenFOAM evaluates $y^+$ at the cell centre, so a correct stack reads half the
nominal figure. Condition~CR is designed at $0.541$ and $0.270$, cruise at $80.3$ and $40.1$.
Measured band medians are $0.279$ to $0.302$ and $37.5$ to $41.9$. Both conditions sit on their
design value, and the cruise wall condition is blended across the buffer layer.

\begin{table}[htbp]
  \centering
  \scriptsize
  \caption{Mesh recipe and as-built statistics. Only the layer stack differs between the two
  conditions. Surface cell size is $1.0/2^{11} = \num{0.488}$~mm. Design $y^+$ is quoted at the
  first-cell top and at the cell centre, the centre being where OpenFOAM's \texttt{yPlus} object
  evaluates. Measured $y^+$ is the band median over the five sampled semispan stations of each
  solved trim case, six at Condition~CR and six at early cruise, from the section cuts. The cruise
  column is the early-cruise mesh; late cruise differs only in its first layer.
  Per-geometry cell counts are in Table~\ref{tab:geom:mesh}.}
  \label{tab:setup:mesh}
  \begin{tabular}{lll}
    \toprule
     & Condition CR (wall resolved) & Cruise (wall modelled) \\
    \midrule
    Background        & \multicolumn{2}{l}{$60\times30\times60$~m at $60\times30\times60$ cells, base $1.0$~m} \\
    Surface level     & \multicolumn{2}{l}{11, giving $\num{0.488}$~mm; feature and TE levels also 11} \\
    Distance shells   & \multicolumn{2}{l}{level 8 / 7 / 6 / 5 at $0.030$ / $0.100$ / $0.350$ / $1.200$~m} \\
    Wake, cone, morph & \multicolumn{2}{l}{level 7} \\
    Quality limits    & \multicolumn{2}{l}{non-orthogonality 65 (relaxed 75), skewness 4} \\
    Decomposition     & \multicolumn{2}{l}{192 subdomains, hierarchical} \\
    \midrule
    Layers asked      & 11, first $4.5$~$\mu$m, $r$ $1.450$ & 5, first $43.7$~$\mu$m, $r$ $1.300$ \\
    Stack asked       & $0.586$~mm & $0.395$~mm \\
    Layers achieved   & $9.79$ of 11, $95.2$\,\% of stack & $4.96$ of 5, $99.5$\,\% of stack \\
    Extrusion         & $97.1$\,\% of wall faces & $99.91$\,\% of wall faces \\
    Wall faces        & $\num{5276772}$ & $\num{5276562}$ to $\num{5277084}$ \\
    Max non-orthogonality & $70.19$ to $70.98$ & $74.10$ to $74.86$ \\
    \midrule
    Design $y^+$, top / centre  & $0.541$ / $0.270$ & $80.3$ / $40.1$ \\
    Measured $y^+$, band median & $0.279$ to $0.302$ & $37.5$ to $41.9$ \\
    \bottomrule
  \end{tabular}
\end{table}

\subsection{The automated pipeline}
\label{sec:setup:pipeline}

\textbf{Pipeline.} Figure~\ref{fig:pipeline} is the route from a delivered geometry to a quotable
number, and every case in this report went down it without manual dictionary editing. The shaded
boxes are gates: a surface that is not closed, a mesh outside the quality limits, or a force
history that has not settled stops the case there. Nothing is asserted by the step that was
supposed to cause it. The run card is written from the case, the harvest is written from the run
cards, and every table and figure in this report is regenerated from the harvest by script, with
the checks of Section~\ref{sec:cond:coverage} run before any number is quoted.

\begin{figure}[htbp]
  \centering
  \resizebox{\linewidth}{!}{\input{fig/pipeline_diagram.tex}}
  \caption{The automated pipeline, one case from geometry to record. Shaded boxes are gates, and
  a case that fails one stops there.}
  \label{fig:pipeline}
\end{figure}
\subsection{Trim convention: every comparison is at matched lift}
\label{sec:cond:trim}
\textbf{Purpose.} Drag is compared at the same lift, not at the same angle of attack. A morphed
wing carries the target lift at a different incidence from the baseline. Comparing the two at one
fixed angle would read a lift difference as a drag difference. This is the single convention the
whole report rests on.

\textbf{Method at Condition~CR.} Three angles bracketing the expected trim are solved to convergence. Three
geometries trim below their first bracket and carry a fourth bracket angle, lower than the other
three. $C_L$ and $C_D$ window means are fitted by least squares through all bracket angles, and
the fit residuals are printed. The trim angle follows from that fit, and a separate trim case is
solved at it: the fourth case run for most geometries, and the fifth for the three that needed an
extra bracket. That solved case is the quotable point, not an interpolation. The target must lie inside the bracket, and an extrapolated
angle is refused.

\textbf{Two details that keep the convention honest.} Incidence is read from each case's own
freestream vector, not from its directory name; \texttt{CMPC\_a0p87} actually runs
$0.8655$~deg. All three bracket cases are asserted to share one operating point, derived from each
case's own density and agreed to $10^{-12}$, before any fit is formed.

\textbf{Method at cruise.} The cruise trims were reached by continuation rather than by
bracketing. Each case converged a first leg at an initial incidence, was nudged in $\alpha$ on its
own lift slope (a prior slope for the first nudge, the secant through its own legs after that),
and continued until $|\Delta C_L| < 10^{-4}$. Each nudge rotates the freestream and the
force-coefficient axes together, and the solved axes are asserted to agree with the trimmed
incidence to $10^{-6}$~deg. Nine of the twelve needed one nudge and three needed two;
Table~\ref{tab:trim:legs} lists every leg. The off-trim legs converged on the same gate as the
trims and are the only points the cruise drag-due-to-lift slopes use.

\begin{table}[htbp]
\centering
\footnotesize
\caption{Every cruise trim leg. $\Delta\alpha$ is the nudge applied after the leg. LCF's third
leg stopped after one iteration and was continued as its fourth, so its legs read 1, 2, 4. Source:
\texttt{data/rans\_forces.json}, \texttt{trim\_history} of each case.}
\label{tab:trim:legs}
\begin{tabular}{lrrrrrrl}
\toprule
Geom. & Leg & Iter. & $\alpha$ (deg) & $C_L$ & $C_L-$tgt (ct) & $C_D$ (ct) & $\Delta\alpha$ (deg) \\
\midrule
\inputrows{sections/gen/tab_trim_legs.tex}
\bottomrule
\end{tabular}
\end{table}

\textbf{Result.} The trim outcome for every geometry and condition is in the performance tables,
Tables~\ref{tab:perf:cr}, \ref{tab:perf:ec} and~\ref{tab:perf:lc}, which carry each trim's
incidence and its lift residual against that condition's own target;
Figure~\ref{fig:cond:clalpha} is the picture of the Condition~CR procedure. The project tolerance
is $|\Delta C_L| < 1\times10^{-4}$, one count of lift, and all eighteen trims are inside it: the
widest is \texttt{LCW\_trim} at $0.895$ counts. Each family's own drag-due-to-lift slope converts
a lift residual into drag, and the largest such bias at any condition is $0.036$ counts, on
\texttt{LCW\_trim}. At cruise six of the twelve slopes rest on a nudge too small to resolve the
slope, under $0.002$ in $C_L$; they are flagged in the data file as fit to the trim residual only,
and they bound it well below $0.1$ counts either way. The same conversion over the six
Condition~CR trims gives at most $0.003$ counts.

\begin{figure}[htbp]
  \centering
  \includegraphics[width=0.40\textwidth]{fig/cl_alpha_condition_CR.pdf}
  \caption{The trim procedure made visible at Condition~CR; the cruise trims were reached by
  continuation and are tabulated in Table~\ref{tab:trim:legs}.
  $C_L$ against $\alpha$ over each family's solved angles, with the target drawn as a dashed
  line and each solved trim case filled. Rings are single-point VSPAERO trims,
  drawn to show agreement on incidence and carrying no lift slope of their own.
  DSO basis, half model, $A_\mathrm{ref} = \num{0.620462}$~m$^2$.}
  \label{fig:cond:clalpha}
\end{figure}

\textbf{Decision.} Every drag count in this report is a difference taken at matched lift, inside
one condition, on one reference frame. The residuals above are worth hundredths of a count of
drag at most, against candidate differences of tenths of a count at Condition~CR and up to fifteen
counts at cruise. The comparison is therefore read as a drag
comparison at fixed lift, with no lift correction applied to any number.

\subsection{Convergence gate}
\label{sec:setup:gate}
\textbf{Method.} A steady SIMPLE solve on this wing limit-cycles about a stationary mean. The
gate is therefore written on the drift of the mean and on the span across the window. It is
evaluated once at the end of each leg, over the last 200 force samples. All three terms must
pass, and the verdict is re-derived from the forces by a second script.
Table~\ref{tab:setup:gate} gives the thresholds and what the 55 cases converged at the
15~September regeneration achieved. \textbf{That set is not the 90 converged cases counted in
Table~\ref{tab:project-scope}}; the rest landed after it. The gate itself is unchanged and every
later case passed it on its own force history. Over the 27 cruise trims and legs that carry every
cruise result in this report, the
achieved values are $C_d$ drift $0.0001$ to $0.0240$~ct (median $0.0022$), $C_l$ drift $0.0037$ to
$0.4007$~ct (median $0.0559$) and $C_l$ span $0.0013$ to $0.8637$~ct (median $0.1223$), all 27
CONVERGED over 200 samples; the widest are first legs, before the nudge, and every trim sits well
inside. What is not claimed is that the tabulated statistics were recomputed over all ninety.

\textbf{Sampling.} 200 samples is a different number of iterations per condition. Condition~CR
writes forces every fifth time step and cruise every step. On a leg ending at \texttt{endTime}
$4000$, which is 46 of those 55, the windows therefore span $t = 3005$ to $4000$
and $t = 3801$ to $4000$. They shift with \texttt{endTime} on the other nine, whose values are
given under Run length below; the window is always the last 200 samples, never a fixed $t$.

\textbf{Reported value.} Every force quoted in this report is a 200-sample window mean. 43
markers carry the gate's own window means. The remaining 12 are recomputed from the merged force
history over the same window. That route reproduces the 43 window-mean markers to
$5.0\times10^{-4}$ counts in both $C_d$ and $C_l$.

\textbf{Run length.} \texttt{endTime} takes three values across the campaign: 46 cases at 4000,
7 at 3000 and 2 at 4500. Every converged case reached its \texttt{endTime}.

\begin{table}[htbp]
  \centering
  \scriptsize
  \caption{The convergence gate and the margins achieved over the 55 cases converged at the
  15~September regeneration. Drift is
  the difference of the two half-window means and span is the first-to-last excursion across the
  window. Drag in counts on the DSO basis, half model,
  $A_\mathrm{ref} = \num{0.620462}$~m$^2$; $C_l$ terms in counts of $C_L$. Values are each case's
  own marker, harvested 2026-09-15T14:05:48Z from
  \texttt{data/rans\_forces\_current.json}.}
  \label{tab:setup:gate}
  \begin{tabular}{llll}
    \toprule
    Term & Threshold & Achieved over the 55 & Median \\
    \midrule
    Window $W$   & 200 samples     & 200 on all 55        & 200 \\
    $C_d$ drift  & $\leq 0.05$~ct  & $0.0000$ to $0.0130$ & $0.0012$ \\
    $C_l$ drift  & $\leq 1.0$~ct   & $0.0004$ to $0.2590$ & $0.0615$ \\
    $C_l$ span   & $\leq 1.0$~ct   & $0.0008$ to $0.5374$ & $0.1068$ \\
    Verdict      & all three terms & CONVERGED, 55 of 55  & n/a \\
    \bottomrule
  \end{tabular}
\end{table}

\subsection{What the evidence covers}
\label{sec:cond:coverage}
\label{sec:project:scope}
\textbf{Coverage.} Ninety cases have converged: 25 at Condition~CR, 33 at early cruise and 32 at
late cruise. The campaign is complete; every geometry is solved at every operating point. Each
passed the same gate on its own force history, over a 200-sample window. Mean drift stays below
$0.05$ counts in $C_D$ and $1.0$ count in $C_L$. The $C_L$ span across the window stays below
$1.0$ count. Every quoted force is that window mean.

\textbf{What the cruise results rest on.} Of the 65 cruise cases, 27 carry every cruise number in
this report: the twelve trims and the fifteen off-trim legs that reached them
(Table~\ref{tab:trim:legs}). The other 38 are alpha brackets kept in the data file and used by no
cruise result.

Table~\ref{tab:project-scope} is the deliverable: six geometries, three operating points, and the
solved cases behind each.

\textbf{Result.} All three conditions are complete at six trims of six. Trimmed lift lands within
$0.10$ counts of target at Condition CR, within $0.57$ counts at early cruise and within $0.90$
counts at late cruise.

\begin{table}[htbp]
  \centering
  \footnotesize
  \caption{Converged cases by geometry and operating point. The three points are
  Condition CR ($M=0.10$, $C_L = 0.428277635108$), early cruise ($M=0.78$,
  $C_L = 0.529297087450$) and late cruise ($M=0.78$, $C_L = 0.544114800210$). Every column
  carries its geometry's solved trim. At cruise, ``trim and legs'' counts the trim and its
  off-trim legs and ``brackets'' the alpha brackets no cruise result uses. Source:
  \texttt{data/rans\_forces.json}, rows generated by \texttt{scripts/make\_cruise\_tables.py}.}
  \label{tab:project-scope}
  \begin{tabular}{llccccc}
    \toprule
     & & Condition CR & \multicolumn{2}{c}{Early cruise} & \multicolumn{2}{c}{Late cruise} \\
    \cmidrule(lr){3-3}\cmidrule(lr){4-5}\cmidrule(lr){6-7}
    ID & Concept & solved & brackets & trim and legs & brackets & trim and legs \\
    \midrule
\inputrows{sections/gen/tab_coverage_rows.tex}
    \bottomrule
  \end{tabular}
\end{table}

\textbf{Two data files, and why both names appear.} \texttt{data/rans\_forces\_current.json} is
the raw harvest: one record per case, read off each case's own marker and \texttt{forceCoeffs}
output. \texttt{data/rans\_forces.json} is the structured merge of that harvest with a per-case
provenance pass, and it is the file that carries the frame assertions, the trim deltas and the
coverage partitions. The performance tables cite the harvest because that is what they were built
from; Figure~\ref{fig:perf:trimmed} cites the merge because that is what it reads. The two
agree on every value quoted in those tables, and that is checked rather than assumed: the figure
writes every plotted number to \texttt{fig/fig\_trimmed\_performance\_values.csv} at full
precision, and all 90 comparisons against the three performance tables, 18 rows by $\alpha$,
$C_D$, $\Delta C_D$, $C_L-$tgt and $L/D$, agree to the precision printed.

\section{Quantities of interest}
\label{sec:qoi}

\textbf{Purpose.} This section is the contents of the results that follow, as a table. It names
every quantity the campaign reports, the frame each one carries, how much of the campaign it
covers, and whether the low-order study produces it at all, which is what decides where a
comparison can be made and where the RANS stands alone.
\subsection{Quantities and their frames}
\label{sec:qv:qoi}
\label{sec:cond:frame}
\textbf{Frame.} Every coefficient here is on the DSO basis. The RANS cases solve a half wing, so
$A_\mathrm{ref} = \num{0.620462}$~m$^2$ and $l_\mathrm{ref} = \num{0.393957}$~m. Twice
$A_\mathrm{ref}$ matches the DSO $S_\mathrm{ref}$ to $3.2\times10^{-6}$ relative. All 90 converged
cases carry that one frame, uniform within each condition. Differences are formed inside a single
condition, never across two.

\textbf{Definitions.} Table~\ref{tab:qv:qoi} is the register. One drag count is $10^{-4}$ in a
force coefficient. Coverage is stated per quantity, because the force set and the surface set
were harvested on different days: forces cover 90 cases and 18 trims, surface data all 18 trims.
The last two columns say whether the low-order study produces the quantity, and where in this
report it is compared.

\begin{table}[htbp]
  \centering
  \scriptsize
  \caption{Quantities of interest, their definitions and their frames. DSO basis throughout:
  ``DSO half'' is $A_\mathrm{ref} = 0.620462$~m$^2$ with $l_\mathrm{ref} = 0.393957$~m, half of
  $S_\mathrm{ref} = 1.24092$~m$^2$ to $3.2\times10^{-6}$. $q$ is the freestream dynamic pressure,
  kinematic at Condition CR and absolute at cruise. Surface quantities carry no reference area.
  ``VLM'' says what the low-order study produces for the quantity and ``Compared'' whether this
  report sets the two side by side. Source: \texttt{data/rans\_forces\_current.json} and \texttt{data/qoi\_methods.md}.}
  \label{tab:qv:qoi}
  \setlength{\tabcolsep}{3pt}
  \begin{tabular}{lllllll}
    \toprule
    Quantity & Definition & Frame & RANS coverage & VLM & Compared & Section \\
    \midrule
    $C_L$, $C_D$ & \texttt{forceCoeffs} on the \texttt{wing} patch & DSO half & 90 cases & $C_{Dtot}$, not compared & no & \ref{sec:perf:cr}, \ref{sec:cruise} \\
    $\Delta C_D$ & candidate minus baseline at fixed $C_L$ & DSO half & 18 trims & induced only & ordering & \ref{sec:qv:counts} \\
    $C_M$ & \texttt{forceCoeffs} about $(1.34,0,0)$~m & DSO half, $l_\mathrm{ref}$ & per case & other frame & no & \\
    $C_{D,p}$, $C_{D,v}$ & \texttt{forces}, same $\hat{d}$ and same $q$ & DSO half & both conditions & no & no & \ref{sec:qv:qoi} \\
    $L/D$ & $C_L/C_D$ at the trim point & area cancels & 18 trims & no & no & \ref{sec:perf:cr}, \ref{sec:cruise} \\
    Trim $\alpha$ & bracket, slope fit, then a solved case & solver \texttt{dragDir} & 18 trims & yes & yes & \ref{sec:qv:lift} \\
    $c_p$ & $(p - p_\infty)/q$, formed in the solver & $q$ per condition & 90 bands & no & no & \ref{sec:app:sections} \\
    $c_f$, $c_{f,s}$ & $|\tau_w|/q$, and its signed form & per-file $q$ & 90 bands & no & no & \ref{sec:app:sections} \\
    $y^+$ & \texttt{yPlus} object and surface field & cell centre & per case & no & no & \ref{sec:setup} \\
    Spanwise load & $c_l c/C_\mathrm{ref}$, pressure only & $C_\mathrm{ref}$ DSO & 18 trims, 78 stations & 69 strips & yes & \ref{sec:qv:span} \\
    $C_{D,i}$, $e$ & Trefftz integral of $(v'^2 + w'^2)$ & half-wing $A_\mathrm{ref}$ & 4 planes, 6 trims & near and far & yes & \ref{sec:cr:spaneff} \\
    $M_h$ & hinge moment about $x_h/c = 0.62$ & local chord & 18 trims & not computed & no & \ref{sec:hinge} \\
    $M_x$ & root bending, solver moment at $y = 0$ & N\,m & 18 trims & yes & yes & \ref{sec:rbm} \\
    \bottomrule
  \end{tabular}
\end{table}

\textbf{Drag split.} The split comes from a \texttt{forces} object on the same drag direction and
the same $q$ as the coefficients. The test is closure against the solver's own $C_D$. On
\texttt{SWB\_a1p90} the parts give $91.218 + 104.098 = 195.316$~ct against 195.317~ct. At cruise
the closure holds once $q$ carries $\rho$: \texttt{CMPB\_a1p20} gives
$131.534 + 51.917 = 183.451$~ct against 183.451~ct.

\textbf{Wall quantities.} $c_p$ is formed inside the solver, so its normalisation travels with the
file. That $q$ is 578~m$^2$/s$^2$ at Condition CR, 11292.8~Pa at early cruise and 8257.37~Pa at late
cruise, each verified
against its own independently known value. Skin friction is carried twice, as a magnitude for
level and as a signed freestream-aligned component for direction. The sign multiplier is measured
on inboard mid-chord faces whose attachment is known, and came out 100~per~cent one sign on all
eighteen trims. Separation is therefore read from a sign rather than a threshold. Median
upper-surface $y^+$ is 0.28 to 0.30 at Condition CR and 37 to 42 at both cruise points.

\textbf{Spanwise load and induced drag.} Sectional load is taken from the wing surface at 78
stations, pressure only. The test is again a closure: the strips integrate back to the case's own
$C_L$ as $C_L = 2\int c_l c\,\mathrm{d}y/S_\mathrm{ref}$, with the explicit half-wing factor of
two. That closure reads 4.65 to 4.91~per~cent low across the six Condition~CR trims and 4.85 to
5.14~per~cent low across the twelve cruise trims, the gap being the root and tip outside the sampled window $\eta = 0.025$ to $0.986$. Induced drag is
defined as a Trefftz integral on four planes per case at $x_{TE} + fc$, $f = 0.05$ to $0.75$. The
induced drag those planes give, and the span efficiency it implies, are reported in
Section~\ref{sec:cr:spaneff}.

\section{Results, quantity by quantity}
\label{sec:results}
\label{sec:perf}

\textbf{Purpose.} Each quantity of Table~\ref{tab:qv:qoi} in turn, with the RANS result first and
the low-order study's value beside it wherever the study produces one. Every coefficient is on the
DSO basis and every difference is formed inside one operating point; the comparison of the two
methods as a selection tool is drawn together in Section~\ref{sec:lo}.

\subsection{Trimmed drag at Condition CR}
\label{sec:cr}
\label{sec:perf:cr}
\label{sec:project:result}

\textbf{Result.} Table~\ref{tab:perf:cr} gives the low-speed result, and the headline is how
little separates the candidates. The five span \SI{1.1}{counts}, from W at $-0.868$ to M at
$+0.232$. Four of the five sit below the baseline and one above it, but every one of them is
inside a count and a half. On $L/D$ the whole field lies between 22.75 and 22.88, a spread of an
eighth of a point; W reaches $22.882$ against the baseline's $22.776$. Every candidate reaches target lift
at lower incidence than the baseline, by 0.19 to 0.50~deg: the aft camber does its intended job
at this condition.
\textbf{Ordering.} The concept ordering differs between the two points, and at low speed only
the grouping is resolvable, not the sequence within it. At Condition CR W and F are the two cheapest, separated by
$0.053$ counts, with F sitting $0.806$ counts below H; at early cruise C and M are the two
cheapest and F is the most expensive of all five. Differences of a twentieth of a count between adjacent pairs at low speed
are the size of the convergence gate itself, so the low-speed sequence is read as a group rather
than a ranking, even though the $1.1$-count spread across the whole field is twenty-two times
that gate and is real. The cruise points, where the field spans forty counts, are where
the ordering carries. Note also that C, M and F are all continuous trailing-edge camber and do
not group together at either point: the concept does not by itself decide the cost.

\begin{table}[htbp]
\centering
\footnotesize
\caption{Trimmed performance at Condition CR, $M=0.10$, target $C_L=0.428277635108$, kOmegaSST
wall resolved. DSO basis, half model, $A_\mathrm{ref}=\num{0.620462}$~m$^2$; one drag count is
$10^{-4}$. $\Delta C_D$ is against \texttt{SWB\_trim} at the same condition and the same target
lift, negative being a saving. $C_L-$tgt is in counts of $C_L$. Rank is by $\Delta C_D$, cheapest
first. Source: \texttt{data/rans\_forces.json}, rows emitted by \texttt{scripts/make\_cr\_table.py}.}
\label{tab:perf:cr}
\begin{tabular}{llrrrrrr}
\toprule
Geom. & Case & $\alpha$ (deg) & $C_D$ (ct) & $\Delta C_D$ (ct) & $C_L-$tgt (ct) & $L/D$ & Rank \\
\midrule
B & SWB\_trim & 1.7476 & 188.031 & $+0.000$ & $-0.094$ & 22.776 & -- \\
C & SWC\_trim & 1.2683 & 187.840 & $-0.191$ & $-0.005$ & 22.800 & 3 \\
M & SWM\_trim & 1.2574 & 188.263 & $+0.232$ & $+0.001$ & 22.749 & 5 \\
F & SWF\_trim & 1.4755 & 187.216 & $-0.815$ & $-0.039$ & 22.876 & 2 \\
W & SWW\_trim & 1.5205 & 187.163 & $-0.868$ & $-0.037$ & 22.882 & 1 \\
H & SWH\_trim & 1.5523 & 188.022 & $-0.009$ & $-0.039$ & 22.778 & 4 \\
\bottomrule
\end{tabular}
\end{table}

\subsection{Trimmed drag at cruise}
\label{sec:cruise}
\textbf{Purpose.} The two cruise points are early cruise at $M = 0.78$ and
$C_L = 0.529297087450$ and late cruise at $M = 0.78$ and $C_L = 0.544114800210$. The two are
separate operating points: each carries its own baseline trim, every increment is formed inside
one of them, and nothing is differenced across them or against Condition~CR. Both run
Spalart--Allmaras wall modelled where Condition~CR runs kOmegaSST wall resolved, so counts from
the two regimes are set beside one another as magnitudes only.

\subsubsection{Early cruise}
\label{sec:perf:ec}
\textbf{Result.} Table~\ref{tab:perf:ec} gives the cruise result, and it is the decision-level
one. \textbf{W is the only candidate that saves drag}, by $0.769$~counts against a baseline of
$177.839$, which raises $L/D$ from $29.764$ to $29.889$. F and H cost $1.796$ and $2.244$~counts,
$0.30$ and $0.37$ points of $L/D$. C and M cost $12.871$ and $14.554$~counts, $2.0$ and $2.3$
points of $L/D$. Every trim lands within $0.565$~ct of target $C_L$, inside the one-count
tolerance; the drag bias that residual implies is at most $0.022$~ct and is carried per case in
the data file.

\textbf{Incidence.} Each candidate again trims at lower incidence than the baseline, by $0.17$
to $0.36$~deg. The commanded camber delivers its lift at cruise as designed. The drag
consequence depends on where it puts the recompression (Section~\ref{sec:perf:mechanism}).

\begin{table}[htbp]
\centering
\footnotesize
\caption{Trimmed performance at early cruise, $M=0.78$, target $C_L=0.529297087450$,
Spalart--Allmaras wall modelled. Frame and conventions as Table~\ref{tab:perf:cr}, with the
baseline \texttt{CMPB\_trim}; rows ranked by $\Delta C_D$.
Source: \texttt{data/rans\_forces.json}, rows generated by
\texttt{scripts/make\_cruise\_tables.py}.}
\label{tab:perf:ec}
\begin{tabular}{llrrrrrr}
\toprule
Geom. & Case & $\alpha$ (deg) & $C_D$ (ct) & $\Delta C_D$ (ct) & $C_L-$tgt (ct) & $L/D$ & Rank \\
\midrule
\inputrows{sections/gen/tab_perf_ec_rows.tex}
\bottomrule
\end{tabular}
\end{table}

\subsubsection{Late cruise}
\label{sec:perf:lc}
\label{sec:perf:late}

\textbf{Result.} Table~\ref{tab:perf:lc} gives the late-cruise result, and it reproduces the
early-cruise picture. W again saves drag, by $0.781$~counts against a baseline of
$187.959$; F and H cost $1.348$ and $1.927$; C and M cost $12.214$ and $13.728$. Every trim lands
within $0.895$~ct of target $C_L$, the widest being \texttt{LCW\_trim}, and incidence again falls
below the baseline's, by $0.16$ to $0.35$~deg.
\begin{table}[htbp]
\centering
\footnotesize
\caption{Trimmed performance at late cruise, $M=0.78$, target $C_L=0.544114800210$,
Spalart--Allmaras wall modelled. Frame and conventions as Table~\ref{tab:perf:ec}, with the
baseline \texttt{LCB\_trim}. Same turbulence model and wall
treatment as early cruise, so these counts are directly comparable with
Table~\ref{tab:perf:ec}. Source as Table~\ref{tab:perf:ec}.}
\label{tab:perf:lc}
\begin{tabular}{llrrrrrr}
\toprule
Geom. & Case & $\alpha$ (deg) & $C_D$ (ct) & $\Delta C_D$ (ct) & $C_L-$tgt (ct) & $L/D$ & Rank \\
\midrule
\inputrows{sections/gen/tab_perf_lc_rows.tex}
\bottomrule
\end{tabular}
\end{table}

\subsubsection{The three conditions together}
\label{sec:perf:inversion}

\textbf{Result.} The two cruise points agree with each other exactly, and they agree with
Condition~CR on everything the low-speed field can resolve. Between early and late cruise
\emph{zero} of the ten candidate pairs change order: both give W, F, H, C, M. At Condition~CR the
sequence is W, F, C, H, M (Table~\ref{tab:perf:cr}). The one pair that differs is C and H, which
at Condition~CR lie $0.182$~counts apart, inside the resolution Section~\ref{sec:perf:cr} gives
neighbours at low speed. \textbf{W is the cheapest geometry at all three operating points}, at
$-0.868$, $-0.769$ and $-0.781$~counts. What changes is the magnitude: the field spans
$1.1$~counts at Condition~CR and $15.3$ at early cruise, almost all of it carried by C and M.

\textbf{What the third condition establishes, which two could not.} A single agreement between
two conditions could be a property of the particular pair of points chosen. A third point that
reproduces the second exactly is evidence that the cruise ordering is a property of the flow
regime rather than of one operating point.
\begin{figure}[p]
\centering
\includegraphics[width=\linewidth]{fig/fig_trimmed_performance.pdf}
\caption{Trimmed performance at matched lift, all six geometries, all three conditions. Panels:
(a) trimmed $C_D$, (b) $L/D$, (c) $\Delta C_D$ at Condition CR, (d) $\Delta C_D$ at early cruise,
(e) $\Delta C_D$ at late cruise, and (f) trim incidence. (d) and (e) share one vertical scale;
Condition CR keeps its own or its deltas would be invisible. Every delta is formed inside one
condition against that condition's own B trim. DSO basis, half model,
$A_\mathrm{ref}=\num{0.620462}$~m$^2$. Source: \texttt{data/rans\_forces.json}.}
\label{fig:perf:trimmed}
\end{figure}

\subsubsection{Why the cruise penalty appears, and why W escapes it}
\label{sec:perf:mechanism}
\textbf{Mechanism.} At $M = 0.78$ the upper surface at $\eta = 0.80$ is supersonic over most of
the chord on every geometry, $144$ to $153$ of $200$ chordwise bins below
$c_p^{*}=-0.4931$ at early cruise. What separates the candidates is how that supersonic region
ends. On B and W it recompresses without a shock: the steepest gradient is
$\mathrm{d}c_p/\mathrm{d}(x/c) = 4.22$ and $4.53$, both at $x/c = 0.8725$, below the shock
threshold of $15$. F and H end it in a moderate shock, $17.1$ at $x/c = 0.7525$ and $19.1$ at
$0.7725$, and stay attached. C and M end it in a strong shock, $45.7$ at $0.7375$ and $50.7$ at
$0.7275$, and separate behind it. \textbf{The drag ordering follows the shock strength}: W within
a count of the baseline, F and H about two counts above it, C and M twelve to fifteen. Late cruise
repeats the pattern at slightly lower strength, $16.1$ and $16.8$ on F and H and $43.6$ and
$38.0$ on C and M. At $M=0.10$ the same camber meets subsonic flow and the field stays inside a
count and a half.

\textbf{Where the load moves.} Figures~\ref{fig:perf:loading} and~\ref{fig:perf:loading:lc} show
the redistribution. Every candidate unloads the inboard wing and loads the outboard wing; the
difference from B crosses zero between $\eta = 0.571$ and $0.652$ and peaks between $0.725$ and
$0.800$, at the same stations to three decimals at both cruise points. The inboard and outboard
integrals cancel, as they must at fixed lift. C and M carry the largest outboard increment,
$+0.075$ and $+0.077$ in $c_l c/C_\mathrm{ref}$ at early cruise, against $+0.046$ for F, H and W,
and that peak sits under the outboard shock.

\textbf{Separation.} Two bands separate in the 30 sampled at each cruise point, and they are the
same two: \texttt{CMPC\_trim} and \texttt{CMPM\_trim} at $\eta = 0.80$
(Figure~\ref{fig:perf:shock}), and their late-cruise counterparts. Reversed flow covers
$23.3$~per~cent of C's upper-surface faces in that band at early cruise, from $x/c = 0.743$ to
$0.992$, and $23.9$~per~cent of M's, from $0.730$ to $0.997$; at late cruise the figures are
$24.1$ and $25.3$~per~cent. Every other cruise band reverses at most $0.019$~per~cent of its
upper-surface faces, and nothing separates at Condition~CR. The cases are wall modelled, so the
sign of this result is firm and its extent carries the precision of the section cuts. The strong
shock and the separation occur together, on the two strong-camber candidates, and they are the
two most expensive. Appendix~\ref{sec:app:sections:ec} gives the station detail.

\begin{figure}[htbp]
\centering
\includegraphics[width=0.46\linewidth]{fig/spanwise_loading_early_cruise.pdf}
\caption{Spanwise loading at early cruise, $M=0.78$, every case trimmed to
$C_L=0.529297087450$. Sectional load $c_l c/C_\mathrm{ref}$ against $\eta$, and the difference
from \texttt{CMPB\_trim}, over 78 stations. Sectional $c_l$ is pressure only. DSO basis,
$C_\mathrm{ref}=\num{0.39396}$~m.}
\label{fig:perf:loading}
\end{figure}

\begin{figure}[htbp]
\centering
\includegraphics[width=0.46\linewidth]{fig/spanwise_loading_late_cruise.pdf}
\caption{Spanwise loading at late cruise, $M=0.78$, $C_L=0.544114800210$, against
\texttt{LCB\_trim}. As Figure~\ref{fig:perf:loading}.}
\label{fig:perf:loading:lc}
\end{figure}

\begin{figure}[htbp]
\centering
\includegraphics[width=0.46\linewidth]{fig/fig_cp_shock_separation_CMPC.pdf}
\caption{One of the two separated bands: \texttt{CMPC\_trim} at early cruise, $\eta=0.80$.
Upper-surface $c_p$, signed freestream-aligned skin friction and reverse-flow fraction. The shock
sits at $x/c=0.7375$ with $\Delta c_p=0.724$ and $c_p=-1.154$ upstream, against
$c_p^{*}=-0.4931$. \texttt{CMPM\_trim} separates at the same station and is drawn beside C in
Figure~\ref{fig:perf:cp:ec}. $c_p$ and $c_f$ carry no reference area.}
\label{fig:perf:shock}
\end{figure}

Section pressure and skin friction at the five spanwise stations, at Condition~CR and both cruise
points, are in Appendix~\ref{sec:app:sections}.

\subsection{Trim angle}
\label{sec:results:trim}
\label{sec:qv:lift}
\textbf{Result.} Trim angle is where the two methods agree best (Table~\ref{tab:qv:lift},
Figure~\ref{fig:qv:trimalpha}). VSPAERO matches the solved RANS angle to within
$0.105^{\circ}$ on all five candidates, and to $0.016^{\circ}$ or better on three of them.
Carried through each family's own fitted lift slope, that is a lift error of at most
2.13~per~cent of the target $C_L$. The baseline B is the outlier at $+0.252^{\circ}$, and its deck
is the only Condition CR run to miss its own target, by $1.53\times10^{-5}$ in $C_L$.

\textbf{Slope.} The RANS side fits $\mathrm{d}C_L/\mathrm{d}\alpha = 0.0860$ to $0.0870$ per
degree over 4 or 5 solved angles per family, with residuals of $1.6\times10^{-4}$ to
$1.1\times10^{-3}$ in $C_l$. The decks are single-point, so the VLM side carries no lift slope and
the comparison is made on trim angle, which needs no derivative.

\begin{table}[htbp]
  \centering
  \footnotesize
  \caption{Trim angle at the common Condition CR target $C_L = 0.428277635108$, both sides on the
  DSO basis, conversion factor exactly 1.000000. RANS $\alpha$ comes from the solved trim case's
  own \texttt{dragDir} header. Implied $\Delta C_L$ is $\Delta\alpha$ through that family's own
  fitted RANS slope, as a percentage of the target. \textbf{Each slope is fitted over that
  family's published alpha-sweep legs together with its Condition~CR trim case, and the trim
  case alone was solved again on the rebuilt mesh: each slope is therefore fitted over legs
  from two different meshes and is not a single-mesh derivative.} They enter only $\Delta\alpha$ converted to
  an implied $\Delta C_L$, which is at most $2.13$ per cent of target, and no drag number in
  this report is formed through them. Rows generated by
  \texttt{scripts/make\_rank\_table.py} from \texttt{fig/fig\_vlmerr\_values.json}.}
  \label{tab:qv:lift}
  \begin{tabular}{lrrrrr}
    \toprule
    Geom. & $\alpha_\mathrm{RANS}$ [deg] & $\alpha_\mathrm{VLM}$ [deg] & $\Delta\alpha$ [deg] & implied $\Delta C_L$ [\%] & $\mathrm{d}C_L/\mathrm{d}\alpha$ [/deg] \\
    \midrule
    B & 1.747609 & 1.99981 & $+0.2522$ & $+5.06$ & 0.085980 \\
    C & 1.268311 & 1.25652 & $-0.0118$ & $-0.24$ & 0.086042 \\
    M & 1.257423 & 1.24972 & $-0.0077$ & $-0.15$ & 0.086119 \\
    F & 1.475456 & 1.45936 & $-0.0161$ & $-0.33$ & 0.087035 \\
    W & 1.520545 & 1.42193 & $-0.0986$ & $-1.99$ & 0.086510 \\
    H & 1.552310 & 1.44697 & $-0.1053$ & $-2.13$ & 0.086784 \\
    \bottomrule
  \end{tabular}
\end{table}

\subsection{Spanwise load}
\label{sec:cr:loads}
\textbf{Result.} Figure~\ref{fig:geom:spanwise} shows the sectional load defined in
Section~\ref{sec:qv:qoi} for every geometry at its trim, with the VSPAERO loading beside it. The
absolute curves lie on top of one another and show the planform. The difference from baseline
isolates the morphing, and it sits over the morphed span in every case; how closely the two
methods agree on that redistribution is Section~\ref{sec:qv:span}.
\begin{figure}[htbp]
  \centering
  \includegraphics[width=0.94\linewidth]{fig/spanwise_loading_condition_CR.pdf}
  \caption{Spanwise loading at Condition~CR, $M = 0.10$, every geometry trimmed to
  $C_L = 0.428277635108$ on the DSO basis. The upper row is the sectional load
  $c_l c / C_\mathrm{ref}$ against $\eta$, where the six geometries differ by a few per~cent and
  the curves overlie. The lower row is each candidate minus the baseline of the same method,
  which is where the morphing appears. The RANS difference peaks at $\eta = 0.7985$ for C, F, H
  and M, and at $\eta = 0.7237$ for W, the one candidate whose command peaks furthest inboard.
  Peak differences are $0.079$ (M), $0.075$ (C), $0.040$ (W), $0.037$ (H) and $0.036$ (F). RANS
  curves carry 78 dense stations over $\eta = 0.025$ to $0.986$; VSPAERO curves carry 69 strips.
  Generated by \texttt{fig/make\_spanwise.py} from \texttt{data/spanwise/} and
  \texttt{data/vlm.json}.}
  \label{fig:geom:spanwise}
\end{figure}

\label{sec:qv:span}

\textbf{Result.} The VLM puts the load where the RANS puts it (Figure~\ref{fig:geom:spanwise}).
Each side is differenced against its own baseline, which is the quantity the optimiser acts on.
Over 78 stations the difference between the two redistributions has an r.m.s.\ of 0.0015 to
0.0037 in $c_l c/C_\mathrm{ref}$, against RANS peaks of 0.034 to 0.080. That is 4.5 to
9.5~per~cent of the peak. The station of peak change agrees to within 0.049 in $\eta$, and to
within 0.014 on four of the five candidates.
\textbf{Resolution.} The RANS gives 78 stations and the VLM 69 strips, interpolated onto a common
grid with an error bound of $8.5\times10^{-4}$. The VLM strips close to their own polar $C_L$ to
$4.6\times10^{-6}$ once the half-wing factor of two is applied, and to 50.0~per~cent without it.
The RANS loading is integrated from the wing surface, pressure only; the VLM loading comes from
the 69 strips of each deck's own \texttt{.lod}. Both are on $C_\mathrm{ref} = 0.39396$~m, the DSO
basis.

\subsection{Induced drag and span efficiency}
\label{sec:cr:spaneff}
\label{sec:spaneff}
\textbf{Method.} Induced drag is taken from the crossflow kinetic energy crossing four
transverse planes in the wake, at $x = 2.394$, $2.626$, $2.916$ and $3.206$~m, sampled by a
function object inside the solver rather than reconstructed afterwards. Span efficiency
follows from it as
%
\begin{equation}
  e \;=\; \frac{C_L^{2}}{\pi \, AR \, C_{Di}} ,
  \label{eq:spaneff}
\end{equation}
%
with $AR$ the full-wing aspect ratio on the DSO basis. Every geometry is evaluated at its
own trimmed angle of attack, and the freestream component is removed using that case's own
$\alpha$; using a different one would leave a uniform residual across the whole plane.

\textbf{Frame.} $C_{Di}$ is on the half wing, $A_\mathrm{ref} = 0.620462$~m$^2$, matching
each case's own force coefficients. $e$ is on the full-wing aspect ratio from the DSO basis,
$b_\mathrm{ref} = 3.6576$~m and $S_\mathrm{ref} = 1.24092$~m$^2$. Condition~CR,
$U_\infty = 34.0$~m/s, every geometry trimmed to $C_L = 0.428277635108$.

\textbf{There is no plateau, so the spread is the uncertainty.} Crossflow kinetic energy is
conserved in an inviscid wake, so a Trefftz integral should not decay downstream at all.
What decay remains is numerical: the grid dissipating the trailing vortex system between one
plane and the next. Quoting a single station would therefore report a number without the
quantity that qualifies it, and Figure~\ref{fig:spaneff} plots all four.

\textbf{The limit at $e = 1$ is a check, not a target.}
For a planar wake $e = 1$ is the elliptical optimum and cannot be exceeded. That makes it
useful twice over: as the reference a spanload is measured against, and as a check on the
calculation itself. A value above~1 would report less induced drag than the theoretical
minimum, which is not a marginal result to be discussed but an arithmetic impossibility, so
any geometry returning one is withheld rather than plotted. Every value in
Figure~\ref{fig:spaneff} sits below the line, as it must.

% RESULT PARAGRAPH: generated. Run
%     python3 fig/make_spaneff.py
%     python3 scripts/spaneff_prose_values.py
% and paste what the second prints. Do NOT hand-write these figures: the decay percentages
% and e values move with every re-trim, and this section is the one place in the report
% where a stale e could be mistaken for a physical result rather than a numerical one.
\textbf{Result.} Span efficiency at the last station ranges from $0.9673$ on the rigid baseline B to $0.9859$ on F, every value below the planar-wake limit. Induced drag at that station spans $54.93$ to $55.98$ counts.

\textbf{Decay between the first and last plane} is $9.5$ to $10.2$ per cent, and is consistent across the geometries: B $9.46$\,\%, W $9.68$\,\%, F $9.95$\,\%, H $9.98$\,\%, C $10.19$\,\%, M $10.22$\,\%. That consistency matters more than the magnitude, because a decay rate that differed between geometries would mean the comparison was being made on wakes resolved to different degrees, and no difference between them could be attributed to the wing.

\textbf{Against the baseline.} Every candidate measured here carries a higher span efficiency than the rigid baseline's $0.9673$ (C, F, H, M and W). The margin is small, one to two per cent, which is the size of spanload improvement the low-order design study predicted for this class of change.


\begin{figure}[htbp]
  \centering
  \includegraphics[width=\linewidth]{fig_spaneff.pdf}
  \caption{Trefftz-plane induced drag and the span efficiency it implies, Condition~CR,
  $M = 0.10$, every geometry trimmed to $C_L = 0.428277635108$. Panel~(a) is $C_{Di}$ at
  each of the four sampled stations; panel~(b) is $e$ from
  Equation~\ref{eq:spaneff}, with the planar-wake limit $e = 1$ marked. $C_{Di}$ is on the
  half wing, $A_\mathrm{ref} = 0.620462$~m$^2$; $e$ is on the DSO-basis full-wing aspect
  ratio. Generated by \texttt{fig/make\_spaneff.py} from
  \texttt{data/span\_efficiency.json}, which is assembled by
  \texttt{scripts/build\_span\_efficiency.py}, which refuses to mix meshes, to mix
  untrimmed legs with trimmed ones, or to carry a value above the physical limit.}
  \label{fig:spaneff}
\end{figure}

\textbf{Span efficiency.} The two methods agree on the level once the far-field definition is
used. VSPAERO reports $E_w = 0.960$ to $0.972$ on the five candidates against its own near-field
$E = 0.568$ to $0.587$, and the RANS Trefftz integral returns $0.972$ to $0.986$ on the
five candidates. The near-field $E$ is not a competing estimate of
the same quantity; it is the one that produced the sign violation reported in
Section~\ref{sec:qv:vlmmethod}.

\subsection{Hinge moment about the morph line}
\label{sec:loads}
\label{sec:hinge}
\textbf{Purpose.} Two structural quantities bear on the concept selection. The hinge moment about
the morph line is the load the actuator reacts, and the low-order method cannot produce it. Root
bending is the active screen the design study selected on, and the low-order method does produce
it, so it is the one structural quantity where the two methods compare directly. Both are reported
here for every geometry at every operating point that carries them.

\textbf{Purpose.} This section reports the aerodynamic moment the actuator has to react. It is
the one quantity in this campaign that the low-order method cannot produce at all: the DSO's own
per-case summary records \texttt{hinge\_moment\_status} as
\texttt{not\_computed\_requires\_surface\_pressure\_integration} for every one of its 23 runs, so
there is no VLM column to compare against and none is constructed here. The decks do carry a
$C_{My}$, but that is the whole-model pitching moment about the deck's own reference point, a
different quantity in a different frame, and reaching a hinge moment from it needs a transfer
formula this project refuses (\texttt{scripts/hinge\_moment.py}).

\textbf{Definition, as agreed with the DSO team.} Every element below is the agreed form and none
of it is a local choice.

\begin{enumerate}
  \item \textbf{Axis:} the swept spanwise line through $x_h/c = 0.62$ at each local section, with
    the moment projected on that line's \emph{local tangent}. The line is swept and tapered, so
    its tangent is not the $y$ axis and projecting on $y$ instead is the error a careful reader
    makes first.
  \item \textbf{Sign:} positive is trailing-edge down. The actuator reaction is the opposite sign.
  \item \textbf{Scope:} surfaces aft of the line only, over $0.60 \leq \eta \leq 1.00$, half wing.
    Aft-of-line is tested per station against $x_h$ at each face's own $y$, never against a single
    global $x$ cut, which on a swept wing would include forward surface at the root and exclude
    aft surface at the tip.
  \item \textbf{Sectional:} $m'_h$ in N\,m/m and $c_{mh} = m'_h/(q_\infty c(y)^2)$ on the
    \emph{local} section chord, because this wing has several defensible reference chords.
  \item \textbf{Integrated:} $M_h$ in N\,m and
    $C_{Mh,\mathrm{region}} = M_h/[q_\infty\!\int\! c(y)^2\,\mathrm{d}y]$ over the same region,
    which avoids selecting one of them.
  \item \textbf{Terms:} pressure is the primary result, the viscous term is reported separately,
    and the total is given.
\end{enumerate}

\textbf{Result.} Table~\ref{tab:hinge} gives every geometry at every condition that carries a
solved trim. \textbf{Every candidate raises the hinge moment above the rigid baseline, at all
three operating points, without exception.} The increase is $2.0$ to $5.0$~per~cent at Condition~CR,
$0.5$ to $6.3$~per~cent at early cruise and $0.4$ to $6.2$~per~cent at late cruise. H is the
largest at every condition, which is what a discrete hinged surface should do: it turns a rigid
segment against the flow rather than redistributing camber smoothly. W is the smallest at both
cruise points.

\textbf{The ordering is preserved between the conditions, up to one near-tie.} From largest
magnitude to smallest it is H, M, C, F, W, B at both cruise points and H, M, C, W, F, B at
Condition~CR, where W and F differ by $0.0006$~N\,m, $0.07$~per~cent of the baseline moment. A
hinge moment measured at low speed therefore ranks the candidates for cruise, as the drag does
(Section~\ref{sec:perf:inversion}).

\textbf{Both orderings rank the candidates on the common comparison line, not on the load each
concept's own actuator carries.} For C, F and M those are the same thing; for H and W they are
not, and W's differs in region as well as in axis. Section~\ref{sec:hinge:hcaveat} states which
rows are actuator moments and which are comparisons only. Read it before using any row to size
a mechanism or to rank the concepts against one another on actuator load.

\textbf{The viscous term is small and it is reported, not assumed.} It runs $0.37$ to
$0.44$~per~cent of the total at Condition~CR and $0.13$ to $0.17$~per~cent at cruise. It is
computed from the sampled wall-shear vector, whose sign convention is measured rather than
inferred: OpenFOAM's \texttt{wallShearStress} is the negative of the traction on the body, and a
wall shear that integrates to a thrust along the freestream is refused outright.

\begin{table}[htbp]
  \centering
  \footnotesize
  \caption{Hinge moment at trim, every geometry and condition, all integrated about the common
  comparison line $x_h/c = 0.62$ over $0.60 \leq \eta \leq 1.00$. $\Delta$ is against that
  condition's own baseline B. \textbf{The last column is load-bearing: only the rows marked
  ``yes'' are the moment that concept's own actuator reacts.} (a) H's hinge is at $x/c = 0.70$,
  so the strip from $0.62$ to $0.70$ lies forward of it and is inside this integral but outside
  H's actuator load. (b) W rotates the whole section about $x/c = 0.25$, so its actuator reacts a
  \emph{torsional} moment over the entire section rather than a hinge moment about $0.62$ over the
  aft $38$ per cent; that quantity is not computed in this campaign. Rows (a) and (b) are valid
  comparisons on a shared line and must not be used to size a mechanism or to rank concepts on
  actuator load. Source: \texttt{fig/hinge\_moment\_values.csv},
  rows generated by \texttt{scripts/make\_cruise\_tables.py}.}
  \label{tab:hinge}
  \begin{tabular}{llrrrrrl}
    \toprule
    & Geometry & $M_h$ [N\,m] & $\Delta M_h$ & $\Delta$ \% & $C_{Mh,\mathrm{region}}$ &
      visc.\ \% & actuator? \\
    \midrule
\inputrows{sections/gen/tab_hinge_rows.tex}
    \bottomrule
  \end{tabular}
\end{table}

\textbf{Reading the two figures.} Both carry one panel per condition, and in both the upper row is
\emph{dimensional} while the lower row is the coefficient. The three conditions sit at
$q = 708.05$, $\num{11292.80}$ and $\num{8257.37}$~Pa, so an $M_h$ or an $m'_h$ from one panel may
not be read against another; $C_{Mh,\mathrm{region}}$ and $c_{mh}$ may, which is why both rows are
drawn. Sectional loading falls monotonically outboard in every case, following the chord, and the
candidates separate from the baseline over the commanded span.

\textbf{The sectional lines are smoothed for display and the integral is not smoothed at all.}
The curves carry a Hann average over $8$~per~cent of the band. What it removes is an artefact of
the reporting bins rather than a feature of the flow: the raw values alternate from station to
station at a lag-one autocorrelation of $-0.49$ on all three baselines, and a block rebin to $91$
stations left that at $-0.49$ while shrinking only the amplitude, which is what established it is
not a physical oscillation. Smoothing moves it to $+0.89$ and takes the direction reversals from
$427$ to $1$, leaving the distribution's shape untouched. The unsmoothed stations are kept in the
per-case records under \texttt{results/derived/}. Nothing is smoothed before integration: $M_h$ in
Table~\ref{tab:hinge} is a direct sum over faces and never passes through these stations.

\begin{figure}[htbp]
  \centering
  \includegraphics[width=\linewidth]{fig_hinge_sections.pdf}
  \caption{Sectional hinge moment against $\eta$. Upper row $m'_h$, lower row $c_{mh}$ on the
  local chord. Lines are smoothed for display.}
  \label{fig:hinge:sections}
\end{figure}

\begin{figure}[htbp]
  \centering
  \includegraphics[width=\linewidth]{fig_hinge_integrated.pdf}
  \caption{Integrated hinge moment, half wing. Upper row $M_h$, lower row
  $C_{Mh,\mathrm{region}}$. The printed value is the total; the viscous share is annotated.}
  \label{fig:hinge:integrated}
\end{figure}

\subsubsection{Which of these numbers are actuator moments, and which are only comparisons}
\label{sec:hinge:hcaveat}

All six geometries are integrated about $0.62$ because the agreed definition fixes one common
comparison line, and $C_{Mh,\mathrm{region}}$ means nothing unless all six share it. That makes
every number in Table~\ref{tab:hinge} a valid \emph{comparison} on a common basis. It does not make
all of them actuator moments, and the distinction is not cosmetic for anyone sizing a mechanism.

\textbf{Three of the five morphed candidates act on the line integrated about.} C, F and M are
continuous trailing-edge camber: the morph starts at $x/c = 0.62$, the surface aft of that line is
the surface that moves, and the moment reported is the moment their actuator reacts.

\textbf{H's own hinge is not at $x_h/c = 0.62$.} It is at $0.70$ (Section~\ref{sec:geom:law}). The
strip between $x/c = 0.62$ and $0.70$ lies \emph{forward} of its real hinge, so it is inside the
integral reported here and outside the load H's own actuator carries. H's number is therefore
correct as a comparison and is not H's actuator moment.

\textbf{W is not a hinged concept at all, and its row is the one most easily misread.} Distributed
twist rotates the whole local section about $x/c = 0.25$, and it is the one candidate whose leading
edge moves (Section~\ref{sec:geom:law}). Its actuator reacts a \emph{torsional} moment, taken about
the quarter chord and over the \emph{entire} section, whereas the quantity tabulated here is taken
about $0.62$ and over the aft $38$~per~cent only. The axis is displaced by $0.37c$ rather than H's
$0.08c$, and the region of integration is wrong as well as the line. \textbf{W's number is a
common-basis comparison and is not, even approximately, W's actuator moment.} The torsional moment
about $x/c = 0.25$ has not been computed in this campaign; it is the same integral over a different
region about a different line, and it is the number the structures and actuator group needs before
distributed twist can be sized or ranked against trailing-edge morphing on actuator load.

\subsubsection{Method uncertainty, measured}
\label{sec:hinge:uncertainty}

\textbf{The spanwise resolution is not free, and it was quantified rather than chosen.} The hinge
line is built from a planform binned in $y$, and the local chord is taken as the extreme $x$ within
each bin, so a coarse bin overstates the chord, moves $x_h = x_\mathrm{LE} + 0.62c$ against the
surface, and changes which faces count as aft of the line. Sweeping the bin count on
\texttt{SWB\_trim} gives $M_h = 0.9375$, $0.8972$, $0.8774$, $0.8677$ and $0.8630$~N\,m at
$100$, $200$, $400$, $800$ and $1600$ bins, with the selected face count falling from
$\num{523627}$ to $\num{501559}$ and the mean chord from $0.2217$ to $0.2135$~m while the mean
$x_h$ holds to $0.015$~per~cent. The increments halve on each doubling, so the convergence is
first order and Richardson puts the limit near $0.858$~N\,m.

\textbf{Every number in this section is at 1600 bins}, which sits about $0.5$~per~cent from that
limit, and $0.5$~per~cent is the figure to quote as the method uncertainty on an absolute $M_h$.
The module's default of 200 bins was $4.5$~per~cent high in magnitude and is superseded. The
\emph{orderings} are unaffected: both condition orderings are identical at 200 and at 1600 bins,
so the ranking was robust to a bias the magnitudes were not.

\textbf{Late cruise.} All eighteen condition-by-geometry cells are integrated, so the third
operating point carries its full six-way comparison. Every candidate raises the magnitude of the
moment about the common line, by $0.4$~per~cent on W to $6.2$~per~cent on H, the same direction
and nearly the same size as at early cruise, and the ordering is identical at the two cruise
points. The actuator-load comparison between these concepts therefore does not depend on which
cruise condition is chosen.

\subsection{Root bending, against the low-order method and against its screen}
\label{sec:rbm}
\textbf{Purpose.} Root bending is one of the two inequality screens the design study selected
on, and the study's own report calls it \emph{the active screen} for the low-speed selection.
Unlike the hinge moment of Section~\ref{sec:hinge}, the low-order method \emph{does} produce it,
so this is the one structural quantity where the two methods can be compared directly. This
section asks two things: do they agree, and does the viscous answer still satisfy the screen.

\textbf{What is compared, and on which basis.} The RANS value is the half-wing moment about the
root chord line, taken from each case's own solver moment output about a centre of rotation
asserted to lie on the symmetry plane. It is not integrated from a surface: root bending needs no
chordwise resolution, so the solver's own $M_x$ is both the more direct route and the one
available for every case. The VLM value is the half-wing root bending in each delivered
optimisation record.

\textbf{Two normalisations have to be stated or the comparison is meaningless.} First, \textbf{root
bending is dimensional}. The VLM decks run at $100$~ft/s and this campaign at $34$~m/s, which is
$q = 569.03$ against $708.05$~Pa, a ratio of $1.24431$. Comparing the two moments without that
factor is wrong by $24$~per~cent before any physics enters. Second, \textbf{the two percentages
are against different baselines}: the VLM differences are against the corrected-interpolation
baseline B$^{*}$, whose own record states it sits $-1.67$~per~cent from the original, while the
RANS differences are against B, the surface that was actually meshed and solved. The percentage
is the primary comparison because it cancels $q$; the dimensional column is reported beside it
precisely because a baseline disagreement would otherwise hide inside the percentages.

\textbf{Result: the two methods agree to better than half a percentage point.}
Table~\ref{tab:rbm} gives both. Every candidate agrees within $0.40$ percentage points, and after
the $q$ ratio the dimensional moments agree to between $0.12$ and $0.88$~per~cent. That is far
closer than the drag comparison of Section~\ref{sec:qv:counts}, and the reason is physical: root
bending is the first moment of the spanwise lift distribution, which a vortex-lattice method
resolves well, and viscosity contributes very little to it. The lift centroid moves outboard with
the morphing exactly as the moment does, from $\eta = 0.3985$ on the baseline to $0.4286$ on M.

\textbf{Result: the viscous calculation breaks the screen on C.} The screen is a $6.8$~per~cent
increase over the rigid baseline. C is the design study's strict optimum and sits at
$+6.796$~per~cent in the VLM, which is $99.94$~per~cent of the screen: the study's own report
says the optimum ``lies immediately below'' it. \textbf{The RANS puts C at $+7.170$~per~cent,
which is over.} M was already over in the low-order method, at $+7.124$, and the RANS agrees and
puts it further over at $+7.519$. F, H and W are comfortably inside on both methods.

\textbf{What that does and does not mean.} The design study states the $6.8$~per~cent figure is
``a study-level bound rather than a certified structural limit'', and that its root bending is
approximated. So this is a flag for the structures and actuator group, not a failure against a
certified allowable. What it does establish is that the constraint which \emph{selected} C is not
satisfied once the same quantity is computed viscously, and that the margin it was selected on,
$0.004$ percentage points, was far smaller than the $0.37$ percentage point difference between the
two methods.

\begin{table}[htbp]
  \centering
  \footnotesize
  \caption{Half-wing root bending at Condition~CR. VLM percentages are against B$^{*}$, RANS
  against B. ``VLM$\times$ratio'' scales the VLM moment by $q_\mathrm{RANS}/q_\mathrm{VLM} =
  1.24431$. Source: \texttt{fig/rootbending\_values.csv}.}
  \label{tab:rbm}
  \begin{tabular}{llrrrrrr}
    \toprule
    & Geometry & VLM \% & RANS \% & diff (pp) & VLM$\times$ratio & RANS [N\,m] & dim.\ \% \\
    \midrule
    B & baseline & $+0.000$ & $+0.000$ & $+0.00$ & $137.805$ & $137.133$ & $+0.49$ \\
    C & \texttt{cte\_i002\_c04}  & $+6.796$ & $\mathbf{+7.170}$ & $+0.37$ & $147.171$ & $146.965$ & $+0.14$ \\
    F & \texttt{cfft\_b02\_c01}  & $+3.755$ & $+3.960$ & $+0.20$ & $142.980$ & $142.563$ & $+0.29$ \\
    H & \texttt{chc\_g02\_c06}   & $+3.769$ & $+3.368$ & $-0.40$ & $142.999$ & $141.752$ & $+0.88$ \\
    M & \texttt{mcv2\_i002\_c01} & $+7.124$ & $\mathbf{+7.519}$ & $+0.40$ & $147.623$ & $147.444$ & $+0.12$ \\
    W & \texttt{cffw\_b01\_c01}  & $+3.010$ & $+3.079$ & $+0.07$ & $141.953$ & $141.355$ & $+0.42$ \\
    \bottomrule
  \end{tabular}
\end{table}

\begin{figure}[htbp]
  \centering
  \includegraphics[width=\linewidth]{fig_rootbending.pdf}
  \caption{Root bending increase against each method's own baseline, with the $6.8$~per~cent
  screen; and the dimensional agreement after the $q$ ratio.}
  \label{fig:rbm}
\end{figure}

\subsubsection{One delivered percentage is stale, and it is the one that matters}
\label{sec:rbm:stale}

\textbf{M's delivered increase is wrong in the file.} \texttt{mcv2\_i002\_c01} records a moment of
$118.64405$~N\,m and an increase of $5.3352$~per~cent. Those two are inconsistent:
$118.64405 / 110.75354 = 1.07124$, an increase of $7.124$~per~cent. The moment is the corrected
value while the percentage was formed against the pre-correction baseline. The design study's own
report states the correct figure in prose, ``the unchanged candidate now increases half-wing root
bending by $7.12$~per~cent, above the illustrative $6.8$~per~cent screen'', so this is a known
defect in the delivered record rather than a new disagreement. \textbf{Every percentage in
Table~\ref{tab:rbm} is recomputed from the two moments rather than read}, and the generator
carries the delivered value alongside and flags the disagreement, so the discrepancy is visible
rather than silently corrected. A reader taking the number straight from the CSV would conclude M
passes the screen by $1.5$ percentage points when it fails it by $0.3$.

\subsubsection{Cruise, and why there is no low-order column}
\label{sec:rbm:cruise}

\textbf{The five candidates do not appear in the study's cruise set at all.} That set contains
optimiser seeds and the rigid baseline, indexed by concept rather than by the geometries meshed
here, so there is no VLM root bending to compare against at either cruise point and none is
constructed. The cruise numbers below are RANS alone.

\textbf{Applying the study's own definition of the cruise limit}, which is the cruise rigid
baseline's own moment at that operating point, so that utilisation is the ratio to it:

\begin{center}\footnotesize
\begin{tabular}{lrrrrrr}
  \toprule
  & B & C & F & H & M & W \\
  \midrule
\inputrows{sections/gen/tab_rbm_cruise.tex}
  \bottomrule
\end{tabular}
\end{center}

\textbf{The limit is not the $6.8$~per~cent screen, and carrying that number forward from
Section~\ref{sec:rbm} inverts every verdict in this table.} The screen applied at Condition~CR
admits a $6.8$~per~cent increase over the rigid baseline. The cruise limit is the study's own and
is stricter: it is the baseline moment itself, so the admissible increase is $0$~per~cent and any
utilisation above $1.000$ is over. The two give opposite answers on the same numbers. On the
$6.8$~per~cent screen every candidate here passes, C and M being the closest at $+6.49$ and
$+6.51$~per~cent at early cruise; on the cruise limit every candidate fails. W carries the
smallest increase at both cruise points, $+3.36$ and $+3.25$~per~cent. Neither bound is a
certified allowable, so the question for the structures and actuator group is which to size
against, and this report does not answer it.

\textbf{All five exceed the cruise limit, and that is consistent with the study rather than
contrary to it.} These are low-speed-selected shapes evaluated at cruise, and the study reports
the same effect from the other direction: transferring its late-cruise optima unchanged to early
cruise takes trailing-edge camber and twist to $105.07$ and $104.02$~per~cent of the limit. A
shape selected at one operating point does not carry its bending margin to another. Late cruise
gives the same ordering as early cruise, C and M highest and W lowest.

\subsection{Surface pressure and skin friction}
\label{sec:results:surface}

The surface distributions behind every integrated quantity above are in the appendices: the
five-station cuts and their overviews in Appendix~\ref{sec:app:sections}, and the whole-wing
difference maps in Appendix~\ref{sec:deltamaps}. Nothing separates at Condition~CR. At both
cruise points the two candidates with the largest commanded camber, C and M, separate behind the
recompression shock at $\eta = 0.80$ and the other four do not (Section~\ref{sec:perf:mechanism}).
The low-order study has no counterpart to any of it.

\section{Can the VLM drive the design study?}
\label{sec:lo}
\label{sec:qv}
\textbf{Purpose.} This section measures VSPAERO against RANS on the same six geometries and
answers whether the vortex-lattice method can keep driving candidate selection. The comparison
is made at Condition~CR, where both sides share one frame and one geometry set, and it asks two
things: does the low-order method put the candidates in the order the RANS puts them, and how
far off is it on the size of the effect.

\subsection{Method}
\label{sec:qv:vlmmethod}
\textbf{Method.} VSPAERO is a vortex-lattice solver, and it ran the same six geometries that were
meshed. Each deck is matched to its RANS case by \texttt{.vsp3} SHA-256, never by name. All nine
Condition CR decks read $S_\mathrm{ref} = 13.3572$~ft$^2$, so the conversion onto the DSO basis is
exactly 1.000000. The package also carries a second baseline deck, \texttt{baseline\_corrected},
which the design study ranks against; every VLM increment in this report is against the deck of
the meshed baseline B, and referencing it to the other deck moves all five candidates by one
constant and changes no ordering.

\textbf{Which columns compare.} VSPAERO writes $C_{Dtot} = C_{Do} + C_{Di}$, and the decks close
that identity to $10^{-12}$. $C_{Do}$ is a flat-plate skin-friction estimate, and it supplies
40.1 to 41.1~per~cent of $C_{Dtot}$ on the six meshed geometries. The decks evaluate it at
$Re = 1\times10^{7}$, against $Re_{C_\mathrm{ref}} = 0.92\times10^{6}$ here on the same chord. The
defensible drag comparison is therefore induced against induced, and $C_{Dtot}$ appears nowhere
beside a RANS $C_D$. Each VLM induced increment below sits beside the RANS total increment at the
same fixed $C_L$; the RANS induced drag from the Trefftz planes is in Section~\ref{sec:cr:spaneff},
and its span efficiency enters Table~\ref{tab:lo:selection} beside the VLM's own.

\textbf{Result.} Table~\ref{tab:qv:vlm} gives the VLM side, and Figure~\ref{fig:qv:decomp}(a)
draws its composition against the RANS total. The two induced columns differ from each other more
than either differs from anything else. Near-field span efficiency runs $E = 0.57$ to $0.59$,
far-field $E_w = 0.94$ to $0.97$, on the same runs. An independent lifting-line calculation on the
spanwise loading returns 56.9~ct, which supports the far-field column to 0.4~per~cent. Both
columns are carried through the table that follows.
\begin{table}[htbp]
  \centering
  \footnotesize
  \caption{VSPAERO at Condition CR, $M = 0.10$, each deck trimmed to $C_L = 0.428277635108$.
  Drag in counts on the DSO basis, deck $S_\mathrm{ref} = 13.3572$~ft$^2$, conversion onto the
  report basis exactly 1.000000. $C_{Di}$ is the near-field surface integration and $C_{Diw}$ the
  far-field wake integral; $E$ and $E_w$ are their span efficiencies. B is the surface the RANS
  meshes. Source: \texttt{data/vlm.json}.}
  \label{tab:qv:vlm}
  \begin{tabular}{llrrrrrrr}
    \toprule
    Geom. & Deck & $\alpha$ [deg] & $C_{Do}$ & $C_{Di}$ & $C_{Dtot}$ & $C_{Diw}$ & $E$ & $E_w$ \\
    \midrule
    B & \texttt{baseline\_wing\_only\_refined} & 1.99981 & 64.060 & 91.676 & 155.736 & 57.649 & 0.5907 & 0.9394 \\
    C & \texttt{cte\_i002\_c04}                & 1.25652 & 64.078 & 92.342 & 156.421 & 56.258 & 0.5865 & 0.9627 \\
    M & \texttt{mcv2\_i002\_c01}               & 1.24972 & 64.084 & 92.617 & 156.701 & 56.441 & 0.5847 & 0.9595 \\
    F & \texttt{cfft\_b02\_c01}                & 1.45936 & 63.894 & 94.465 & 158.359 & 55.712 & 0.5733 & 0.9721 \\
    W & \texttt{cffw\_b01\_c01}                & 1.42193 & 63.927 & 95.030 & 158.957 & 55.827 & 0.5699 & 0.9701 \\
    H & \texttt{chc\_g02\_c06}                 & 1.44697 & 63.887 & 95.370 & 159.257 & 56.383 & 0.5679 & 0.9605 \\
    \bottomrule
  \end{tabular}
\end{table}

\begin{figure}[htbp]
  \centering
  \begin{subfigure}[t]{0.40\textwidth}
    \centering
    \includegraphics[width=\textwidth]{fig/drag_decomposition_rans_vs_vlm_condition_CR.pdf}
    \caption{Drag decomposition.}
    \label{fig:qv:drag}
  \end{subfigure}
  \hfill
  \begin{subfigure}[t]{0.40\textwidth}
    \centering
    \includegraphics[width=\textwidth]{fig/trim_alpha_rans_vs_vlm_condition_CR.pdf}
    \caption{Trim angle.}
    \label{fig:qv:trimalpha}
  \end{subfigure}
  \caption{Condition CR, $M = 0.10$, $C_L = 0.428277635108$, DSO basis, RANS half model
  $A_\mathrm{ref} = 0.620462$~m$^2$. (a) What each method contains: the RANS total beside the
  VSPAERO $C_{Dtot} = C_{Do} + C_{Di}$ split into its parasite correlation and its induced part,
  with the far-field $C_{Diw}$ alongside. The two bar heights are not differenced, because
  $C_{Do}$ is a flat-plate correlation at the deck Reynolds number. (b) Trim angle at the common
  target $C_L$, which is the closest agreement in this section.}
  \label{fig:qv:decomp}
\end{figure}

\subsection{Drag increments}
\label{sec:qv:counts}
\textbf{Result.} Table~\ref{tab:lo:selection} is the central comparison of this section. The
RANS morphing signal spans $-0.868$ to $+0.232$~ct, $1.100$ in total, and neither VLM column
resolves an effect that small. The near-field column misses the measured increment by $0.709$ to
$4.222$~counts and gets the sign wrong on all four candidates the RANS puts below the baseline.
The far-field column misses by $0.955$ to $1.441$~counts, every error of the same sign and within
half a count of the others, which is why its ordering survives where the near-field ordering does
not: a near-constant offset moves the whole field together and leaves the sequence intact.

\textbf{Far field.} On the far-field column every predicted increment is negative. The VLM has all
five candidates saving induced drag, while the RANS measures four of the five costing total drag.
Those errors run 1.24 to 7.85~ct, and for F, W and H they exceed the entire signal span. Part of
that gap is a real profile increment the VLM has no mechanism to produce, and a RANS Trefftz
integral is what separates the two.

\textbf{Systematics.} The RANS-side systematic uncertainty is far smaller than any of these
differences. Window drift is
0.0014 to 0.0028~ct against a gate of 0.05~ct, and removing each trim case's own $C_L$ residual
moves the increments by at most 0.021~ct.

\begin{table}[htbp]
  \centering
  \scriptsize
  \setlength{\tabcolsep}{3.5pt}
  \caption{Selection test at Condition CR, $M = 0.10$, $C_L = 0.428277635108$, DSO basis (RANS
  half model $A_\mathrm{ref} = 0.620462$~m$^2$; VSPAERO $S_\mathrm{ref} = 13.3572$~ft$^2$,
  conversion exactly 1.000000). Increments are against the baseline B, in counts: the RANS
  increment is a total drag increment and the VLM increments are induced drag increments, so
  ``error'' is the VLM prediction minus the RANS measurement. Ranks are by increasing increment
  against each method's own baseline, rank 1 best; the closest adjacent RANS pair is W and F at $0.053$~ct, so no pair is tied. The last
  three columns are span efficiencies: the VLM's own near- and far-field values, and the RANS
  value at the last Trefftz station. Rows and statistics generated by
  \texttt{scripts/make\_rank\_table.py} from \texttt{fig/fig\_vlmerr\_values.json},
  \texttt{data/vlm.json} and \texttt{data/span\_efficiency.json}.}
  \label{tab:lo:selection}
  \begin{tabular}{lrrrrrrrrrrr}
    \toprule
     & \multicolumn{2}{c}{RANS} & \multicolumn{3}{c}{VLM near field} &
       \multicolumn{3}{c}{VLM far field} & \multicolumn{3}{c}{span efficiency} \\
    \cmidrule(lr){2-3}\cmidrule(lr){4-6}\cmidrule(lr){7-9}\cmidrule(lr){10-12}
    Geom. & $\Delta C_D$ & rank & $\Delta C_{Di}$ & rank & error & $\Delta C_{Diw}$ & rank & error &
      VLM $E$ & VLM $E_w$ & RANS $e$ \\
    \midrule
    W & $-0.868$ & 1 & $+3.354$ & 4 & $+4.222$ & $-1.823$ & 2 & $-0.955$ & 0.5699 & 0.9701 & 0.9837 \\
    F & $-0.815$ & 2 & $+2.789$ & 3 & $+3.604$ & $-1.938$ & 1 & $-1.123$ & 0.5733 & 0.9721 & 0.9859 \\
    C & $-0.191$ & 3 & $+0.666$ & 1 & $+0.857$ & $-1.392$ & 3 & $-1.201$ & 0.5865 & 0.9627 & 0.9773 \\
    H & $-0.009$ & 4 & $+3.694$ & 5 & $+3.703$ & $-1.266$ & 4 & $-1.257$ & 0.5679 & 0.9605 & 0.9773 \\
    M & $+0.232$ & 5 & $+0.941$ & 2 & $+0.709$ & $-1.209$ & 5 & $-1.441$ & 0.5847 & 0.9595 & 0.9724 \\
    \midrule
    \multicolumn{3}{l}{Kendall $\tau$, all five}  & \multicolumn{3}{c}{$-0.20$} & \multicolumn{3}{c}{$+0.80$} & & & \\
    \multicolumn{3}{l}{Spearman $\rho$, all five} & \multicolumn{3}{c}{$-0.20$} & \multicolumn{3}{c}{$+0.90$} & & & \\
    \multicolumn{3}{l}{Pairs swapped, of 10}      & \multicolumn{3}{c}{6} & \multicolumn{3}{c}{1} & & & \\
    \bottomrule
  \end{tabular}
\end{table}

\subsection{Is the ranking preserved?}
\label{sec:qv:rank}
\label{sec:lo:rank}
\label{sec:project:ranking}
\textbf{Result.} Ranking is the property a search tool needs, and it survives on one column
(Table~\ref{tab:lo:selection}), but not the column the June selection was made on. The RANS orders the
candidates $W < F < C < H < M$. Far-field $C_{Diw}$ gives $F < W < C < H < M$, nine of ten pairs
concordant, Kendall $\tau = +0.80$. Near-field $C_{Di}$ gives $C < M < F < W < H$, four of ten
concordant, $\tau = -0.20$.

\textbf{Which way the two metrics fail.} The far-field metric's single error is on the closest
pair in the field: the RANS separates W and F by $0.053$~counts, the narrowest adjacent gap of the
five, and against a convergence gate of $0.05$~counts per window that pair is at the edge of what
this campaign can resolve. A metric that orders everything except the one pair nobody can
separate is performing about as well as the data permits. The near-field metric fails differently
and not by a small margin: $\tau = -0.20$ is worse than no information, because it groups
$\{C, M\}$ as the cheap pair and $\{F, W, H\}$ as the expensive one, and the viscous answer is
close to the reverse.

\textbf{What this claim is not.} The Condition~CR field spans $1.100$~counts and the convergence
gate is $0.05$~counts per window, so the spread is twenty-two gates wide but the closest adjacent
pair is one. Reproducing this ordering is therefore evidence that a metric tracks the sign and the
grouping of the field, not that it resolves neighbours within it; Section~\ref{sec:perf:cr}
makes the same point from the other side, that the low-speed sequence is read as a group. The two
statements are consistent and the distinction is the whole of it: a search tool that puts the
cheap candidates in the cheap group is useful even where it cannot separate two of them, and that
is what $\tau = +0.80$ with its one error on the narrowest pair describes.
\textbf{Reading.} The two induced-drag columns the VLM writes are not interchangeable, and this
is where the choice between them shows. The far-field wake integral is the one that tracks the
viscous ordering; the near-field surface integration does not.

\begin{figure}[htbp]
  \centering
  \includegraphics[width=0.74\linewidth]{fig_vlmerr_ranking.pdf}
  \caption{Selection test at Condition CR, $M=0.10$, every geometry trimmed to
  $C_L = 0.428277635108$ on the DSO basis. Each panel ranks the five candidates by drag
  increment, low order on the left and RANS on the right, rank 1 being the least drag. Panel (a)
  uses near-field $C_{Di}$, the metric the June selection was made on, and swaps 6 of 10 pairs.
  Panel (b) uses far-field $C_{Diw}$, the metric the design study now selects on, and swaps 1 of
  10, that pair being W against F at $0.053$~counts apart. Generated by
  \texttt{fig/make\_vlmerr.py} from \texttt{data/vlm.json} and the Condition CR RANS trims.}
  \label{fig:project-ranking}
\end{figure}

\subsection{Decision}
\label{sec:qv:decision}
\textbf{Decision.} The VLM can be trusted to order candidates, on its far-field column only, and
to price none of them. Ordering: far-field $C_{Diw}$ reproduces the RANS order to nine pairs of
ten, near-field $C_{Di}$ to four of ten. Pricing: the whole morphing effect at Condition CR spans
$1.100$~counts, from W at $-0.868$ to M at $+0.232$, and the far-field increments miss the
measured ones by $0.955$ to $1.441$~counts, which is $87$ to $131$~per~cent of that span. The
near-field increments miss by $0.709$ to $4.222$~counts, $64$ to $384$~per~cent.

\textbf{Why one of them still ranks.} The five far-field errors are all of one sign and lie
within half a count of each other, $-0.955$ to $-1.449$. An error that is nearly a constant offset
moves every candidate together and therefore leaves the ordering intact while destroying the
magnitude, which is precisely the behaviour seen. The near-field errors are neither
one-signed in magnitude nor clustered, spanning a factor of six, and the ordering does not
survive them.

\textbf{Consequence.} Candidate search can stay with VSPAERO provided it selects on the
far-field column, which is the column the design study adopted in its August correction and not
the one the June selection was made on. Every count quoted for a selected concept comes from
RANS, and at Condition CR that is not a refinement but a requirement: the low-order error is the
size of the entire effect being resolved.

\textbf{Consequence for selection.} Both metrics predict a saving for every candidate at
Condition~CR and the viscous calculation returns a saving for four of the five. At cruise it
returns a saving for W alone and a cost for the other four (Section~\ref{sec:cruise}). \textbf{The
three geometries the low-order study currently selects, W, F and H, are the three cheapest at both
cruise points, in that order}. C and M, from the July selection on the near-field column, cost
twelve to fifteen counts.
The cruise cost of C and M is wave drag and separation behind a strong aft shock
(Section~\ref{sec:perf:mechanism}), which no induced-drag column can see, so the selection is
right at cruise for a reason the low-order method does not model. Reselection that admits
strong-camber shapes should therefore rank on a metric that sees the shock; of the two low-order
induced-drag columns, the far-field one is the only one that orders the candidates at low speed.

\section{What this report establishes, and what it does not}
\label{sec:concl}

\textbf{Purpose.} This section states the whole result in one place, after the loads that bear on
it, so that no single section can be read as the conclusion on its own.

\subsection{What is established}
\label{sec:concl:established}

\begin{enumerate}
  \item \textbf{Drag: W is the cheapest geometry at every operating point, and the low-speed
    ordering carries to cruise.} At Condition~CR four candidates improve on the rigid baseline,
    W by $0.868$, F by $0.815$, C by $0.191$ and H by $0.009$ counts against $188.031$, and the
    whole field spans $1.1$ counts, enough to separate the groups but not the neighbours within
    them. At early cruise W saves $0.769$ counts against $177.839$, F and H cost $1.796$ and
    $2.244$, and C and M cost $12.871$ and $14.554$; late cruise reproduces it, W at $-0.781$ to
    M at $+13.728$ against $187.959$. \textbf{The two cruise points agree exactly}, both giving W,
    F, H, C, M, and they differ from the low-speed sequence W, F, C, H, M only in the C and H pair,
    which is $0.182$ counts apart at Condition~CR (Section~\ref{sec:perf:inversion}).
  \item \textbf{The cruise cost is the aft shock.} C and M end the supersonic region at
    $\eta = 0.80$ in a strong shock and separate behind it; F and H in a moderate shock and stay
    attached; W and the baseline without a shock. Drag follows that order exactly
    (Section~\ref{sec:perf:mechanism}).
  \item \textbf{Hinge moment: every candidate raises the actuator load, and two of the rows are
    not actuator loads at all.} All five raise the moment about the common $x_h/c = 0.62$ line at
    all three operating points, by $2.0$ to $5.0$ per cent at Condition~CR, $0.5$ to $6.3$ at
    early cruise and $0.4$ to $6.2$ at late cruise, with the same ordering at all three up to one
    near-tie. C, F and M act on the line integrated about, so their numbers are the
    moment their actuator reacts. \textbf{H's and W's are not}: H's hinge is at $0.70$, and W is
    not hinged at all but rotates the whole section about $0.25c$
    (Section~\ref{sec:hinge:hcaveat}).
  \item \textbf{Root bending: the constraint that selected C is not satisfied viscously.} C sits
    at $+6.796$ per cent in the VLM, inside the study's $6.8$ per cent screen by $0.004$
    percentage points, and the RANS puts it at $+7.170$ per cent, outside. M is outside on both.
    F, H and W are comfortably inside on both. At both cruise points, under the study's own
    stricter limit, all five exceed the baseline moment, W by the least
    (Section~\ref{sec:rbm:cruise}).
  \item \textbf{The two methods agree on root bending to better than half a point.} VLM and RANS
    differ by at most $0.40$ percentage points, so the screen result above is a property of the
    shapes rather than of the method.
  \item \textbf{The low-order method orders the candidates on its far-field column and prices
    none of them.} Far-field $C_{Diw}$ reproduces the RANS order at Condition~CR to nine pairs of
    ten, its one error on the closest pair in the field; near-field $C_{Di}$, the metric the June
    selection was made on, gives four of ten. Neither resolves the size of the effect: the
    far-field increments miss the measured ones by $86$ to $131$~per~cent of the whole low-speed
    span and the near-field by $63$ to $381$ (Section~\ref{sec:qv:decision}).
\end{enumerate}

\subsection{What is not established, stated as plainly}
\label{sec:concl:not}
\label{sec:perf:limits}

\begin{enumerate}
  \item \textbf{W's actuator load is not computed anywhere in this report.} Distributed twist
    reacts a torsional moment about $x/c = 0.25$ taken over the entire section. The quantity
    tabulated for W is about $0.62$ over the aft $38$ per cent, which is a valid comparison on a
    shared line and is not that load. Until the torsional moment is computed, \textbf{trailing-edge
    morphing and distributed twist cannot be ranked against one another on actuator load}, which is
    the axis the concept selection was expected to turn on.
  \item \textbf{H's actuator load is likewise not its tabulated number}, for the smaller reason
    that its own hinge is $0.08c$ aft of the comparison line.
  \item \textbf{Neither root-bending bound is a certified allowable.} The design study calls the
    $6.8$ per cent figure ``a study-level bound rather than a certified structural limit''. The
    cruise limit is a different and stricter definition. Which to size against is a question for
    the structures and actuator group and is not answered here.
  \item \textbf{The low-order comparison is made at Condition~CR only.} The VSPAERO decks and the
    meshed geometries share one frame and one geometry set there and nowhere else: the cruise
    decks carry the TP-1580 $S_\mathrm{ref}$ and a geometry set that maps to none of the meshed
    cases, and the study's cruise set contains optimiser seeds and the rigid baseline rather than
    the five candidates. Extending the comparison to either cruise point requires the
    corresponding VLM columns at that operating point, not merely the RANS ones, and quoting the
    two together without them would cross operating points (D076).
\item The meshed baseline is the uncorrected aerofoil-interpolation lineage and all five
candidates are the corrected one. Measured directly between the two \texttt{.vsp3} files
by \texttt{scripts/measure\_lineage\_pair.py} into \texttt{data/lineage\_measured.json},
the uncorrected wing carries 9 distinct aerofoils across its 19
stations where the corrected wing carries 19, so the two surfaces agree exactly at every
station they share and diverge only in between. They are identical at and inboard of
$\eta = 0.450$; the first station that differs is $\eta = 0.515$, which is also the worst,
at $0.407$ percentage points of $t/c$ and $0.209$ per cent of chord in surface ordinate. The
camber difference is an order of magnitude smaller again, at most $0.0034$ per cent of chord.
Both wings are the same planform at the same twist, so this is a thickness-channel
difference, and the candidate-to-candidate comparison is taken between surfaces of the
corrected lineage throughout and does not cross it.
\item Condition CR runs kOmegaSST wall resolved and both cruise points run Spalart--Allmaras wall
modelled, so counts are differenced within a condition only.
\item Every case is fully turbulent from the leading edge and in free air, so a tunnel comparison
needs blockage, wall-interference and support corrections on the measured side.
\item The VSPAERO columns are induced-drag increments and the RANS columns are total-drag
increments, so Table~\ref{tab:lo:selection} compares ordering rather than magnitude.
\item The cruise drag-due-to-lift slopes are formed on the trim legs alone, and six of the twelve
rest on a nudge too small to resolve the slope. They are used only to convert trim residuals of
at most $0.895$ counts of $C_L$ into a drag bias of at most $0.036$ counts.
\end{enumerate}

\subsection{Recommendation}
\label{sec:perf:recommendation}

\textbf{Recommendation.} W is the geometry to carry forward. It is the cheapest of the six at all
three operating points, $0.868$, $0.769$ and $0.781$ counts below the baseline at Condition~CR,
early cruise and late cruise, it raises $L/D$ at both cruise points, and it carries no aft shock
and no separation at cruise. It also has the smallest root-bending increase of the five at both
cruise points, $+3.36$ and $+3.25$ per cent, and sits inside the $6.8$ per cent screen at
Condition~CR on both methods. F and H are the next two, within $2.3$ counts of the baseline at
cruise. C and M are not candidates for cruise: they cost twelve to fifteen counts, separate behind
a strong aft shock, and both fail the root-bending screen once it is computed viscously.

\textbf{The caveat on W, and it is the one this report cannot close.} W's actuator reacts a
torsional moment about the quarter chord over the whole section, and that moment is not computed
here (Section~\ref{sec:hinge:hcaveat}); its tabulated hinge moment is a common-basis comparison
only. W's aerodynamic case is complete. Its mechanical case is not, and it is the input the
structures and actuator group needs before W can be committed to.

\subsection{The one number that would change the picture}
\label{sec:concl:next}
W's torsional moment about $x/c = 0.25$ over the full section, at every operating point. It is
the same integral this report already performs, over a different region about a different line, and
it is the missing input to the concept selection, and the recommendation above rests on it. Nothing
else outstanding has comparable leverage: the drag case is closed at all three operating points,
and the root-bending case is closed against both bounds.


\appendix
\section{Surface pressure and skin friction, by section}
\label{sec:app:sections}

\textbf{Purpose.} This appendix carries the surface distributions behind the integrated
coefficients in the body of the report. Every geometry is cut at five spanwise stations,
$\eta = 0.20$, $0.40$, $0.60$, $0.80$ and $0.95$, at its own trimmed incidence to that
condition's own target lift: $C_L = 0.428277635108$ at Condition~CR, $C_L = 0.529297087450$ at
early cruise and $C_L = 0.544114800210$ at late cruise, all on the DSO basis. Comparing the six geometries' $-c_p$ peaks at a station
means something only because they share that target within a condition.

\textbf{Conditions are never overlaid.} Condition~CR and cruise are normalised on different
dynamic pressures: $q = 578$ (kinematic, $\mathrm{m^2/s^2}$, the incompressible solver carries no
density) against $q = \num{11292.8}$~Pa at early cruise and $\num{8257.37}$~Pa at late cruise.
Each is correct and they are not comparable curves, so every condition is plotted as its own set
and the plotting script refuses to mix them.

\textbf{Coverage.} All three conditions are six geometries of six: $90$ station files over $18$
cases, with no case partially represented.

\textbf{What is plotted.} Figures~\ref{fig:perf:cp}, \ref{fig:perf:cp:ec} and~\ref{fig:perf:cp:lc} carry $c_p$ on an
inverted axis and its difference from the baseline by surface; Figures~\ref{fig:perf:cf}
and~\ref{fig:perf:cf:lower} carry the signed skin friction by surface. All are against $x/c$ on
the local chord of that station, which falls from $0.4816$~m at $\eta = 0.20$ to $0.1671$~m at
$\eta = 0.95$.

\textbf{The curves are a band, not a line.} The pressure and skin-friction
curves pool faces within $\pm 0.004$ semispan, $\pm 7.32$~mm, and the wing is swept and tapered,
so faces at one $x/c$ sit at different $y$ and carry genuinely different $z$; curves are the median
over $0.005c$ bins per surface rather than the raw scatter. Within one bin the $10$--$90$ percentile spread of $z/c$ is
$1.5$ to $2.3$~per~cent of chord at the leading edge, where the surface turns through nearly
$180^\circ$, and about $1.1$~per~cent at the blunt trailing edge, against $0.21$ to
$0.49$~per~cent at mid-chord. A defective surface would be ragged at mid-chord too.

\textbf{Two of the three sources of raggedness are the sampling band. The third was a defect and
is now corrected.} The leading-edge scatter is the band meeting high curvature: across the five
stations the within-bin spread grows by a factor of $2.79$ while the band's share of the local
chord grows by $2.80$, with mid-chord flat throughout as a control. That is also why the outboard
stations are the noisiest, the fixed $14.63$~mm slab spanning $3.0$~per~cent of chord at
$\eta = 0.20$ but $8.8$~per~cent at $\eta = 0.95$. The inboard trailing-edge ripple is the band
meeting taper rather than curvature: the chord shortens by $0.2279$~m per metre of span, so it
varies by $3.33$~mm across the slab and a face is binned at an $x/c$ that is not quite its own.
Neither is removable without narrowing the band, and both are sampling rather than flow. The
OUTBOARD trailing-edge ripple was neither, and it was not flow either: it was the upper-lower
split failing where the two surfaces converge. Splitting on the exported facet normals instead
removes it. Measured as the within-bin $C_p$ spread over $x/c = 0.92$ to $1.00$, C at
$\eta = 0.80$ falls from $0.2623$ to $0.0204$ and M from $0.2797$ to $0.0177$; at $\eta = 0.95$
the six cases fall from about $0.34$ to about $0.09$. Two controls hold and are what make this a
correction rather than a smoothing: mid-chord is unchanged to four decimal places, $0.0447$
against $0.0447$ for C at $\eta = 0.80$, and the inboard station is unchanged at $\eta = 0.20$,
ratio $1.03$, because the band and the taper are still there.

\textbf{The band also over-reads the local chord, by a fixed length rather than a fixed
fraction.} It takes its leading edge from the inboard face of the slab and its trailing edge from
the outboard one, so on a swept wing it over-reads by
$(w/2)(\tan\Lambda_{\mathrm{LE}} + \tan\Lambda_{\mathrm{TE}})$. With $w = 14.63$~mm,
$\Lambda_{\mathrm{LE}} = 28.84^\circ$ and $\mathrm{d}c/\mathrm{d}y = -0.266$ that predicts
$6.11$~mm against $6.67$~mm measured, and the measured excess is near-constant across span while
the chord itself halves, which is the signature of a slab width rather than of a scaling error.
The chords quoted above are the planar-cut values and are accordingly $1.4$~per~cent smaller than
the band's at $\eta = 0.20$ and $3.8$~per~cent smaller at $\eta = 0.95$. Because $x/c$ is
normalised to span $0$ to $1$, the curves are unaffected; only the dimensional chord printed
above each column of the figures differs.

\textbf{The sign on $C_{f,s}$ is the useful part.} $C_{f,s}$ is the component of wall shear
along the freestream, taken from each case's own \texttt{dragDir}, so \emph{negative means
reversed flow}. The sign convention is calibrated per case against inboard mid-chord faces whose
attachment is known independently, not assumed; it came out one-signed on all twelve cases.

\subsection{Overview at the five stations}
\label{sec:app:overview}

Figures~\ref{fig:perf:cp} to~\ref{fig:perf:cf:lower} draw the five stations of each condition on
one page each, pressure first and then skin friction by surface; the two subsections that follow
read them condition by condition.
\begin{figure}[htbp]
\centering
\includegraphics[width=\linewidth]{fig/fig_cp_sections_CR.pdf}
\caption{Surface pressure at five spanwise stations, $\eta = 0.20$ to $0.95$, Condition~CR, every
geometry at its own trim to $C_L = 0.428277635108$. The top row is $c_p$ on each case's own
freestream $q$, inverted, with the grey band the baseline's interquartile spread within the pooled
span band; the two lower rows are each candidate minus the baseline on the upper and lower
surfaces, drawn from $x/c = 0.03$ because forward of that the branch split and the span band are
ill-conditioned. $x/c$ is on the band's local chord, printed above each column.}
\label{fig:perf:cp}
\end{figure}

\begin{figure}[htbp]
\centering
\includegraphics[width=\linewidth]{fig/fig_cp_sections_earlycruise.pdf}
\caption{Surface pressure at the same five stations, early cruise, every geometry at its own trim
to $C_L = 0.529297087450$; layout as Figure~\ref{fig:perf:cp}. The dashed line is the sonic
$c_p^{*}$, the tick marks at $\eta = 0.80$ locate each candidate's shock by its $x/c$, and the
shaded runs are where the upper-surface flow on C and M is reversed.}
\label{fig:perf:cp:ec}
\end{figure}

\begin{figure}[htbp]
\centering
\includegraphics[width=\linewidth]{fig/fig_cp_sections_latecruise.pdf}
\caption{Surface pressure at the same five stations, late cruise, every geometry at its own trim
to $C_L = 0.544114800210$; layout as Figure~\ref{fig:perf:cp:ec}.}
\label{fig:perf:cp:lc}
\end{figure}
\textbf{Reading the skin-friction figures.} Condition~CR is wall resolved at $y^+ \approx 1$ with
kOmegaSST and cruise is wall modelled with Spalart-Allmaras, so $c_f$ is not equally resolved
between the two conditions and the two columns of each figure are read separately rather than
against one another. Neither $c_p$ nor $c_f$ carries a reference area.

\textbf{The plotted quantity and its sign.} The figures show
$c_{f,s} = m\,(\tau_w\!\cdot\!\hat{d})/q$, the \emph{freestream-aligned} component of wall shear,
with $\hat{d}$ taken from each case's own \texttt{forceCoeffs} drag direction. On this swept wing
that is not the section-normal component. OpenFOAM returns a negative wall shear for flow attached
in $+x$, so a calibrated sign multiplier $m$ is read from every file's own header, $m = -1$ on all
of them, and positive then plots as attached. It is therefore the \textbf{sign change} that
locates separation, never the magnitude. Lines are medians over $0.005c$ bins of the sampled band,
with at least five faces per bin. A band counts as a \emph{separated run} only when at least four
consecutive bins are negative, which is more than $2$~per~cent of local chord, so an isolated
negative bin is not read as separation. At $\eta = 0.95$ and $x/c > 0.875$ the curve is drawn in
grey because the aft bin there is band-averaged rather than resolved.

\textbf{What the ninety bands contain.} Over the $90$ (case, station) bands of the three
conditions, the largest reversed fraction on any upper surface is $25.3$~per~cent
(\texttt{LCM\_trim}, $\eta = 0.80$), and $31$ bands carry any reversed upper-surface face at all.
On the lower surface the largest is $0.474$~per~cent (\texttt{SWC\_trim}, $\eta = 0.20$); $56$
bands carry a trace, and those sit at the root junction and in the blunt trailing-edge base region
rather than being boundary-layer separation. The four bands that genuinely separate, C and M at
$\eta = 0.80$ at each cruise point, are named in Section~\ref{sec:perf:mechanism}.
\begin{figure}[htbp]
\centering
\includegraphics[width=\linewidth]{fig/fig_cf_sections_upper.pdf}
\caption{Signed skin friction $c_{f,s}$ at the same five stations, \textbf{upper} surface,
Condition~CR wall resolved on the left and early cruise wall modelled on the right, each geometry
at its own trim. Positive is attached and negative reversed, so it is the sign change and not the
magnitude that locates separation: the two separated runs at early cruise, C and M at
$\eta = 0.80$, are marked with their separation and reattachment stations. Late cruise separates
at the same two bands (Section~\ref{sec:app:sections:lc}). The
grey region at $\eta = 0.95$ is where the aft bin is band-averaged rather than resolved.}
\label{fig:perf:cf}
\end{figure}

\begin{figure}[htbp]
\centering
\includegraphics[width=\linewidth]{fig/fig_cf_sections_lower.pdf}
\caption{The same, \textbf{lower} surface. No lower-surface band carries a separated run at
either condition; the traces at the blunt trailing edge are the base region, not boundary-layer
separation.}
\label{fig:perf:cf:lower}
\end{figure}
\subsection{Condition CR, \texorpdfstring{$M = 0.10$, $C_L = 0.428277635108$}{M 0.10, CL 0.428277635108}}
\label{sec:app:sections:cr}

Nothing separates. Across all six geometries and all five stations the reversed fraction never
exceeds $0.25$~per~cent of faces in a band, and those faces sit at the blunt trailing edge and
the root junction, where they are present on the rigid baseline too. The six geometries lie
almost on top of one another: at $\eta = 0.60$ the upper-surface suction plateaus agree within
about $0.02$ in $-C_p$, and the morphing appears only aft of the hinge.

\subsection{Early cruise, \texorpdfstring{$M = 0.78$, $C_L = 0.529297087450$}{M 0.78, CL 0.529297087450}}
\label{sec:app:sections:ec}

At $M = 0.78$ the aft-camber candidates carry a recompression shock on the upper surface between
$x/c \approx 0.65$ and $0.78$ that is absent from the baseline. Two of them separate behind it,
and they are the two with the largest commanded camber.

At $\eta = 0.80$, C peaks at $-C_p = 1.279$ at $x/c = 0.716$ and M at $1.249$ at $x/c = 0.703$.
Both lose the boundary layer behind the shock and do not recover it before the trailing edge:
$23.3$~per~cent of C's upper-surface faces in that band are reversed, over $x/c = 0.743$ to
$0.992$, and $23.9$~per~cent of M's, over $x/c = 0.730$ to $0.997$, while both lower surfaces
reverse nowhere at all. M's shock is the stronger of the two, $\mathrm{d}c_p/\mathrm{d}(x/c) =
50.7$ against C's $45.7$, which follows its larger command: $A/c = 0.0350$ against C's $0.0337$. F
and H carry a weaker shock at a lower peak, $0.988$ at $x/c = 0.721$ and $1.027$ at
$x/c = 0.734$, and stay attached with zero reversed faces; F's command is $A/c = 0.0139$, two and
a half times smaller. W peaks forward, $1.074$ at $x/c = 0.067$, and has no aft shock at all. The
baseline peaks at $0.906$ at $x/c = 0.054$ and likewise has none.

This is the one place in the campaign where a candidate's flow is qualitatively different from
the baseline's rather than quantitatively so, and it is what makes C and M the two most expensive
candidates at cruise. It is local in span: at $\eta = 0.60$ and $0.40$ neither C nor M reverses a
single face, and at $\eta = 0.95$ each reverses one face, $0.010$~per~cent.

\subsection{Late cruise, \texorpdfstring{$M = 0.78$, $C_L = 0.544114800210$}{M 0.78, CL 0.544114800210}}
\label{sec:app:sections:lc}

Late cruise repeats early cruise station by station. At $\eta = 0.80$ C and M carry the strong
shock, $43.6$ at $x/c = 0.7325$ and $38.0$ at $0.7175$, peak at $-C_p = 1.248$ and $1.228$, and
separate behind it: $24.1$~per~cent of C's upper-surface faces reversed, over $x/c = 0.734$ to
$0.994$, and $25.3$~per~cent of M's, over $0.725$ to $0.997$, with both lower surfaces clean. F
and H carry the weaker shock, $16.1$ and $16.8$, and stay attached. W and the baseline have no aft
shock. Off $\eta = 0.80$ the picture is the early-cruise one: no reversal on C or M at
$\eta = 0.40$ and $0.60$, and one face each at $\eta = 0.95$.

\textbf{Provenance.} Sections are cut from each case's own \texttt{wingSurface.vtk} at its
converged iteration by \texttt{scripts/section\_slices\_from\_surface.py}, which also writes the
upper/lower flag from the per-face Newell normal, its sign calibrated per case against the more
negative mean $C_p$ at mid-chord and cross-checked against median $z$, refusing rather than
guessing if the two criteria disagree. The wind axes are the trimmed leg's own, read from its
\texttt{forceCoeffs} history in numeric order of the legs. The cruise surfaces were pulled from the
cluster with every file's checksum verified, and each section manifest records the source case,
the surface checksum and the extractor version. The figures are drawn by \texttt{fig/make\_cp.py}
and \texttt{fig/make\_cf.py}. The station files, their per-file headers ($q$, $p_\infty$, local
chord, \texttt{dragDir}, shear sign) and the extraction manifests are in
\texttt{data/sections\_welded/} for both cruise points, \texttt{data/sections\_r2/} for B, C, F,
H and W at Condition~CR and \texttt{data/sections\_nflag/} for M at Condition~CR, resolved per
case by \texttt{fig/section\_source.py}; the superseded set carrying the geometric split is
retained beside them in \texttt{data/sections/}.

% Figures: scripts/render_delta_maps_of12.py --all, from results/wing_maps_of12/, promoted
% into fig/. Six of the nine rendered maps are used here: Cp and the SIGNED Cf for each
% condition. The |Cf| magnitude maps are rendered but NOT carried, because a magnitude cannot
% express reversed flow and the signed field is strictly more informative on the one question
% these figures exist to answer.
\section{Surface differences on the wing}
\label{sec:deltamaps}

\textbf{Purpose.} Section~\ref{sec:app:sections} cuts the wing at five stations. This section
gives the same comparison over the whole surface at once: the baseline field, and each
candidate minus the baseline, on the upper and lower surfaces together. Twelve panels per
quantity per condition, so a feature can be located in span and chord at the same time rather
than inferred from five slices.

\textbf{A difference between two separately meshed wings needs a common support.} The wings
are meshed independently and share no cells at all: SWB carries $11{,}334{,}640$ surface
triangles against SWC's $11{,}356{,}180$, and no face of one corresponds to any face of the
other. A cell-to-cell subtraction therefore does not exist. Every face is instead binned onto a
common $(\eta, x/c)$ grid of $240 \times 320$ cells, area-weighted, with upper and lower
surfaces binned separately using the flag each export carries. That grid is the only support on
which a difference between these two wings is defined, and $76{,}797$ of its cells are filled
on both surfaces.

\textbf{The panels are drawn on the planform, not on the grid.} An $(\eta, x/c)$ image is a
rectangle, and a rectangle silently asserts a rectangular wing to any reader who does not stop
to read the axes. Each bin therefore carries its physical position and the panels are drawn on
the true swept and tapered planform, matching Figure~\ref{fig:geom:displacement} so the two can be read
against one another. Span runs left to right with the root at the left, $x$ increases upward, so
the leading edge is at the bottom. The drawn planform reproduces the registered wing to
$0.18$~per~cent in root chord ($0.6403$~m against $0.639125$), $0.2437$ against $0.238$ in
taper, and $1.0036$~m of leading-edge sweep. It is a canvas for a field, not a measurement of
the planform.

\textbf{What the figures carry.} Column~B is the baseline's own absolute field on its own scale.
The five candidate columns are candidate minus baseline on ONE shared scale, set at the
$99.5$th percentile of $|\Delta|$, so their magnitudes are comparable by eye across columns.
Blank cells are cells no face landed in, and are never zero. Every panel title carries that
case's $C_D$ in counts, its $\Delta C_D$ against the baseline of the same condition, and its
$L/D$, read from the same \texttt{rans\_forces\_current.json} harvest as the tables in
Sections~\ref{sec:cr} and~\ref{sec:cruise}, so a figure and a table cannot disagree. Each wing is
at its own trimmed incidence at the common target $C_L$, so every comparison here is at MATCHED
LIFT.
\textbf{A difference forward of the hinge is lineage, not morphing.} The candidates are not
identical to the baseline ahead of the commanded region. Measured on the placed surfaces in
absolute coordinates, C, F, H and M all depart from the baseline by the same amount, $0.316$,
$0.184$ and $0.079$~mm at $\eta = 0.60$, $0.80$ and $0.95$, and by exactly zero at
$\eta = 0.40$ and inboard; W departs further and differently, $0.366$, $1.171$ and $0.292$~mm.
Four candidates sharing one forward shape that the baseline does not is the baseline-lineage
difference of Section~\ref{sec:geom:lineage} appearing in the meshed surfaces, and it does not
cancel in a
baseline-to-candidate difference because the two sides sit on different lineages. It is small,
$2.6$~per~cent of the aft displacement at $\eta = 0.80$, but any faint forward feature in the
candidate columns should be read as lineage rather than as an effect of the morph.

\textbf{Null.} Differencing the baseline against itself returns exactly zero in every filled
cell, for every binned quantity and all three conditions, $76{,}797$ cells each. The renderer also
refuses to difference across operating points: it compares condition, Mach, $U_\infty$, grid
size and registered semispan between every pair before subtracting, and stops rather than
differencing unlike states.

\subsection{Pressure}
\label{sec:deltamaps:cp}

At Condition~CR the candidates are almost free. The pressure differences are confined to the
commanded region and are small. The trimmed drag behind each panel is Table~\ref{tab:perf:cr},
and the differences there are small beside the cruise ones but they are not scatter: the
systematics this report quantifies, in Section~\ref{sec:qv:counts}, total at most
$0.024$~counts, being window drift of $0.0014$ to $0.0028$~ct and a trim-residual effect of at
most $0.021$~ct. Section~\ref{sec:perf:recommendation} states what a margin of this size does and
does not support.
At cruise the same geometries separate into three groups. The baseline runs $177.8$~counts at
$L/D\ 29.76$ at early cruise; W saves $0.8$, F and H cost $1.8$ and $2.2$, and C and M cost
$12.9$ and $14.6$. Late cruise gives the same order. The aft-camber candidates carry a
recompression shock on the upper surface that the baseline does not, and it shows in these maps
as a sharp spanwise band of strongly negative $\Delta C_p$ running outboard across the rear of
the wing: strongest on C and M, weaker on F and H, and absent on W.

\begin{figure}[htbp]
  \centering
  \includegraphics[width=\linewidth]{fig/delta_map_CR_cp.png}
  \caption{$C_p$ on the wing at Condition~CR, $M = 0.10$: the baseline's own field and each
  candidate minus the baseline, upper and lower surfaces. Each wing at its own trimmed
  incidence at the common target $C_L$, so this is at matched lift.}
  \label{fig:delta:cr:cp}
\end{figure}

\begin{figure}[htbp]
  \centering
  \includegraphics[width=\linewidth]{fig/delta_map_EC_cp.png}
  \caption{$C_p$ on the wing at early cruise, $M = 0.78$. The band of strongly negative
  $\Delta C_p$ across the rear of C, F, H and M is the recompression shock the baseline does not
  carry; W has none. The two most expensive candidates on drag, C at $+12.9$~ct and M at
  $+14.6$~ct, are also the two whose boundary layer does not survive that shock; see
  Section~\ref{sec:deltamaps:cfs}.}
  \label{fig:delta:ec:cp}
\end{figure}

\begin{figure}[htbp]
  \centering
  \includegraphics[width=\linewidth]{fig/delta_map_LC_cp.png}
  \caption{$C_p$ on the wing at late cruise, $M = 0.78$; as Figure~\ref{fig:delta:ec:cp}.}
  \label{fig:delta:lc:cp}
\end{figure}

\subsection{Skin friction and separation}
\label{sec:deltamaps:cfs}

\textbf{The signed component, not the magnitude.} These maps carry $C_{f,s}$, the component of
wall shear along the freestream taken from each case's own \texttt{dragDir}, so negative means
reversed flow. The magnitude $|\tau|/q$ is positive by construction and cannot express
reversal: a separated region reads as merely small and is indistinguishable from a thinned but
firmly attached boundary layer. That distinction is the whole question here, because F carries
the same shock as C and M with a visibly weakened boundary layer behind it and does not reverse.

\textbf{The sign is calibrated, never assumed.} Taken raw, the first case reported
$99.18$~per~cent of its faces reversed, which is not a separated wing but an inverted
convention: OpenFOAM's \texttt{wallShearStress} opposes the freestream where the boundary layer
is healthy. Rather than hard-code a minus sign, the sign is measured per case against faces
whose attachment is known independently of this field, namely inboard mid-chord upper faces at
$\eta = 0.15$ to $0.35$ and $x/c = 0.35$ to $0.65$. It came out one-signed on all eighteen trims,
with about $260{,}000$ calibration faces each.

\textbf{Nothing separates at Condition CR.} All six geometries reverse at most $0.013$~per~cent
of their upper-surface bins, which is $0$ to $10$ bins of $76{,}797$ at the root junction and the
blunt trailing edge, and is present on the rigid baseline too. The six are indistinguishable from
one another and from the noise floor.

\textbf{At both cruise points two candidates separate and four do not.} Measured over the whole
upper surface at early cruise: C reverses $1.520$~per~cent of its bins and M $2.122$, against
$0.004$ for the baseline, F and H and $0.000$ for W. At late cruise C reverses $1.803$ and M
$2.335$, the other four the same as at early cruise. That is a factor of some $400$ between the
two groups, on the same grid and the same code path. The lower surfaces reverse
$0.000$~per~cent on every geometry.

\textbf{The maps agree with the section cuts, on a different support.} At $\eta = 0.80$ at early
cruise the binned maps give $23.28$~per~cent of bins reversed for M and $22.97$ for C, against
the face-level section extraction's $23.86$ and $23.27$~per~cent at the same station; at late
cruise $25.16$ and $23.91$ against $25.28$ and $24.06$. Both methods return exactly zero for B,
F, H and W at both cruise points. Two independent code paths, one binning onto a common
$240 \times 320$ grid and the other counting raw faces in a band, agree to within $0.6$
percentage points and agree exactly on which cases separate. The whole-wing figures above
are smaller only because the whole wing is a much larger denominator: the reversal is real and
confined in span, not damped by the binning.

\textbf{Reduced skin friction is not separation, and the contour is what separates the two.} A
difference panel shows that the shear FELL, not that it REVERSED. The black contour is each
case's own $C_{f,s} = 0$, so flow inside it is reversed and the contour's ABSENCE is itself a
result. It is drawn only where at least $0.1$~per~cent of a surface's filled bins are negative,
because every case carries a handful of negative bins at the root junction and the trailing
edge; that threshold sits about eight times above the worst noise ($10$ bins) and fifteen times
below the smallest real signal ($1167$ bins, for C at early cruise). At Condition CR no panel
draws a contour. At both cruise points only C and M do.

\begin{figure}[htbp]
  \centering
  \includegraphics[width=\linewidth]{fig/delta_map_CR_cf_s.png}
  \caption{$C_{f,s}$ at Condition~CR, $M = 0.10$. No separation contour appears on any panel:
  every geometry is attached everywhere, and the candidates differ from the baseline only in the
  magnitude of the shear over the morphed region.}
  \label{fig:delta:cr:cfs}
\end{figure}

\begin{figure}[htbp]
  \centering
  \includegraphics[width=\linewidth]{fig/delta_map_EC_cf_s.png}
  \caption{$C_{f,s}$ at early cruise, $M = 0.78$. C and M carry a closed separation contour over
  the rear outboard wing; F, H and W do not, although F and H show a fall in shear behind their
  weaker shock. The two most expensive candidates on drag, C at $+12.9$~ct and M at $+14.6$~ct,
  are the two whose boundary layer does not survive it.}
  \label{fig:delta:ec:cfs}
\end{figure}

\begin{figure}[htbp]
  \centering
  \includegraphics[width=\linewidth]{fig/delta_map_LC_cf_s.png}
  \caption{$C_{f,s}$ at late cruise, $M = 0.78$; as Figure~\ref{fig:delta:ec:cfs}.}
  \label{fig:delta:lc:cfs}
\end{figure}


{\footnotesize
\begin{thebibliography}{9}
\setlength{\itemsep}{1pt}
\bibitem{tp1580} Morgan, H.~L., and Paulson, J.~W., \emph{Low-Speed Aerodynamic Performance of
a High-Aspect-Ratio Supercritical-Wing Transport Model Equipped With Full-Span Slat and
Part-Span Double-Slotted Flaps}, NASA TP-1580 (L-13201), December 1979. NTRS 19820007142.
\bibitem{tm80048} Morgan, H.~L., \emph{Model Geometry Description and Pressure Distribution
Data From Tests of EET High-Lift Research Model}, NASA TM-80048, February 1979.
\bibitem{agard138} \emph{Experimental Data Base for Computer Program Assessment}, AGARD
AR-138, 1979. Onera M6 wing, run 308; the validation case for the cruise numerics used here.
\bibitem{zheng} Zheng, L., \emph{ARGUS Aerodynamic Optimization of a Morphing Outer Wing},
internal ARGUS design-space-optimisation report, July 2026, with the far-field $C_{Diw}$ audit
delivery of 19 August 2026. Table~18 is the cruise operating-point table referred to in the
caption of Table~\ref{tab:cond:points}.
\bibitem{openfoam} \emph{OpenFOAM}, version 12, OpenFOAM Foundation. All cases run
\texttt{foamRun} with the \texttt{incompressibleFluid} and \texttt{fluid} modules.
\end{thebibliography}
}

\end{document}
